% ROZDZIAL 15 — Korutyny, zadania i porownanie ze std::async (C++20).
% Fragment do wspolnego skryptu: bez preambuly i bez \begin{document}.
% Wymagane pakiety: listings, amsmath, amssymb, booktabs, tabularx,
% array, hyperref (jak w rozdzialach poprzednich).
% Dolaczenie: % ROZDZIAL 15 — Korutyny, zadania i porownanie ze std::async (C++20).
% Fragment do wspolnego skryptu: bez preambuly i bez \begin{document}.
% Wymagane pakiety: listings, amsmath, amssymb, booktabs, tabularx,
% array, hyperref (jak w rozdzialach poprzednich).
% Dolaczenie: % ROZDZIAL 15 — Korutyny, zadania i porownanie ze std::async (C++20).
% Fragment do wspolnego skryptu: bez preambuly i bez \begin{document}.
% Wymagane pakiety: listings, amsmath, amssymb, booktabs, tabularx,
% array, hyperref (jak w rozdzialach poprzednich).
% Dolaczenie: % ROZDZIAL 15 — Korutyny, zadania i porownanie ze std::async (C++20).
% Fragment do wspolnego skryptu: bez preambuly i bez \begin{document}.
% Wymagane pakiety: listings, amsmath, amssymb, booktabs, tabularx,
% array, hyperref (jak w rozdzialach poprzednich).
% Dolaczenie: \input{rozdzial_15_korutyny_vs_async_cpp.tex}
% Listing oznaczony jako „kompletny program” mozna skompilowac oddzielnie.

\chapter{Korutyny, zadania asynchroniczne i \texttt{std::async}}
\label{chap:cpp15-coroutines}

\section{Po co C++ potrzebuje korutyn?}

Zwykła funkcja ma zasadniczo jeden punkt wejścia. Może wywoływać inne
funkcje, ale po powrocie oddaje sterowanie wywołującemu i nie posiada
mechanizmu późniejszego wznowienia dokładnie w miejscu, w którym
zakończyła wykonywanie. \emph{Korutyna} (ang. \emph{coroutine}) jest
funkcją, której wykonanie można \emph{zawiesić}, a następnie
\emph{wznowić}, zachowując wybrane zmienne lokalne i miejsce kontynuacji.

Korutyny C++20 wspierają między innymi:
\begin{itemize}
\item \emph{generatory} --- wytwarzanie kolejnych wartości na żądanie,
      bez budowania całej kolekcji;
\item \emph{zadania asynchroniczne} --- opis operacji, których kolejne
      etapy uruchamiane są po zakończeniu wcześniejszych;
\item \emph{kooperacyjne planowanie} --- oddawanie sterowania
      harmonogramowi zadań;
\item \emph{maszyny stanów} --- zachowywanie stanu pomiędzy wznowieniami
      bez ręcznego kodowania każdej możliwej kontynuacji.
\end{itemize}

Najważniejsza zasada brzmi:
\begin{center}
\emph{Korutyna nie jest wątkiem, a \texttt{co\_await} nie oznacza
automatycznie ,,uruchom w nowym wątku''.}
\end{center}

Korutyna określa \emph{jak zatrzymać i kontynuować obliczenie};
\texttt{std::async}, \texttt{std::thread} i pula wątków dotyczą również
\emph{gdzie i kiedy} wykonywać pracę. W szczególności kod korutyny
może działać od początku do końca w jednym wątku, nawet jeśli
wielokrotnie się zawiesza. Odwrotnie, korutyna może być wznawiana
przez różne wątki, jeżeli tak zdecyduje dostarczony mechanizm planowania.

\subsection{Trzy słowa kluczowe}

\begin{center}\small
\begin{tabularx}{\textwidth}{@{}p{0.32\textwidth}X@{}}
\toprule
Konstrukcja & Działanie w korutynie \\
\midrule
\code{co_await expr} & Przekazuje sterowanie obiektowi oczekującemu; może zawiesić korutynę bez blokowania wątku. \\
\code{co_yield value} & Udostępnia kolejną wartość i zwykle zawiesza generator do wznowienia. \\
\code{co_return value} & Kończy korutynę i przekazuje wynik zgodnie z protokołem jej typu zwracanego. \\
\code{await_ready()} & Określa, czy wynik jest już gotowy i można pominąć zawieszenie. \\
\code{await_suspend(h)} & Obsługuje zawieszenie i uchwyt \code{h}; może zaplanować późniejsze wznowienie. \\
\code{await_resume()} & Zwraca rezultat oczekiwania albo zgłasza zapisany wyjątek po wznowieniu. \\
\bottomrule
\end{tabularx}
\end{center}

C++20 wprowadził następujące konstrukcje składniowe:
\begin{description}
\item[\texttt{co\_await expr}] Zastosuj protokół oczekiwania do wyrażenia;
      zależnie od niego możesz zawiesić korutynę i później ją wznowić.
\item[\texttt{co\_yield expr}] Udostępnij kolejną wartość generatora;
      składniowo jest to szczególna postać oczekiwania, wykorzystująca
      metodę \texttt{yield\_value} obiektu \texttt{promise\_type}.
\item[\texttt{co\_return expr}] Zakończ korutynę i przekaż wynik
      przez odpowiednie \texttt{return\_value} albo \texttt{return\_void}.
\end{description}

Funkcja zawierająca jeden z tych operatorów (w odpowiednim kontekście)
rozpoznawana jest jako korutyna. Typ jej zwracanej wartości musi
udostępniać właściwy \emph{typ obietnicy} \texttt{promise\_type}
(lub odpowiednią specjalizację \texttt{std::coroutine\_traits}).
To wymóg mechanizmu korutyn, a nie znanego już \texttt{std::promise}.

\begin{lstlisting}[caption={Niepoprawna intuicja: std::future nie jest automatycznie typem korutyny (C++20); fragment celowo błędny.}]
#include <future>

// Standardowe std::future<int> nie zapewnia promise_type:
// std::future<int> calculate() { co_return 42; } // BLAD!
\end{lstlisting}

W C++20 biblioteka standardowa dostarcza \emph{mechanizm} korutyn
(nagłówek \texttt{<coroutine>}), ale nie dostarcza gotowego, ogólnego
\texttt{Task<T>} ani standardowego executora. Takie typy definiuje
biblioteka lub programista. \texttt{std::generator} jest rozwiązaniem
późniejszym, wprowadzonym dopiero w C++23, a jego wsparcie zależy
od użytej biblioteki standardowej.

\section{Różnica między zawieszeniem a zablokowaniem wątku}
\label{sec:cpp15-suspend-block}

Rozważmy dwa sposoby wstrzymania działania programu:
\begin{itemize}
\item \texttt{std::this\_thread::sleep\_for(100ms)} \emph{blokuje
      bieżący wątek} na pewien czas; sam wątek pozostaje zajętym
      zasobem systemowym, choć może nie zużywać wówczas czasu CPU;
\item dobrze zaprojektowane \texttt{co\_await} \emph{zawiesza
      korutynę}, oddając wątek do wykonywania innych zadań;
      wznowienie może nastąpić po nadejściu wyniku operacji I/O,
      komunikatu lub upływie czasu.
\end{itemize}

Nie każde \texttt{co\_await} daje zawieszenie: awaiter może zgłosić,
że wynik jest już gotowy. Również samo użycie \texttt{co\_await}
nie gwarantuje nieblokującego I/O --- zależy to od jego implementacji.

\begin{center}
\small
\begin{tabularx}{\textwidth}{@{}lXX@{}}
\toprule
\textbf{Właściwość} & \textbf{Blokada wątku} & \textbf{Zawieszenie korutyny} \\
\midrule
Zachowanie stanu & Na stosie wywołań wątku & Zwykle w ramce korutyny \\
Wątek potrzebny podczas oczekiwania & Tak, istnieje i czeka & Nie jest konieczny \emph{dla samego oczekiwania} \\
Wznowienie & Planista systemu/odblokowanie & Mechanizm awaitera i właściciel kontynuacji \\
Koszt & M.in. wątek i jego stos & M.in. ramka i obsługa zawieszenia \\
\bottomrule
\end{tabularx}
\end{center}

\section{Ramka korutyny i jej cykl życia}
\label{sec:cpp15-frame}

Kompilator przekształca korutynę w obiekt reprezentujący jej stan.
Ten stan nazywamy \emph{ramką korutyny} (ang. \emph{coroutine frame}).
Może ona obejmować obiekt \texttt{promise\_type}, parametry i stan
potrzebny do kontynuowania obliczenia po zawieszeniu. To, które
zmienne lokalne trzeba zachować, wynika z przebiegu programu.

W ogólnym modelu wywołanie korutyny przebiega w następujących etapach:
\begin{enumerate}
\item Alokowana jest ramka, chyba że implementacja legalnie
      wyeliminuje tę alokację (np. przez zagnieżdżenie ramki w innym obiekcie).
\item Kopiowane lub przenoszone są parametry odpowiednio do ich rodzajów;
      \emph{parametry referencyjne pozostają referencjami}.
\item Konstruowany jest obiekt obietnicy \texttt{promise\_type}.
\item Wywoływana jest metoda \texttt{get\_return\_object()}.
\item Wykonywany jest \texttt{co\_await promise.initial\_suspend()}.
\item Uruchamiane jest ciało korutyny (od razu lub po wznowieniu).
\item \texttt{co\_return} uruchamia \texttt{return\_value} lub
      \texttt{return\_void}; nieobsłużony wyjątek trafia do
      \texttt{unhandled\_exception()}.
\item Niszczone są obiekty lokalne, a korutyna wykonuje końcowy
      \texttt{co\_await} na wyniku \texttt{promise.final\_suspend()}.
\item Ramka zostaje ostatecznie zniszczona we właściwym momencie.
\end{enumerate}

Kolejność i czas \emph{zwolnienia ramki} zależą od przyjętej polityki
własności. Dla \texttt{final\_suspend = suspend\_always} ramkę zwykle
niszczy właściciel uchwytu, nawet po zakończeniu ciała korutyny.
Przy \texttt{final\_suspend = suspend\_never} mechanizm korutyny
może zakończyć się bez pozostawiania ramki zawieszonej do późniejszego
zniszczenia przez właściciela. Tych strategii nie wolno mieszać.

\subsection{Początek i koniec: \texttt{initial\_suspend} i \texttt{final\_suspend}}

Dwa najczęściej używane typy z nagłówka \texttt{<coroutine>} to:
\begin{description}
\item[\texttt{std::suspend\_never}] Oczekiwanie jest natychmiast
      gotowe i nie powoduje zawieszenia.
\item[\texttt{std::suspend\_always}] Oczekiwanie zawsze prowadzi
      do zawieszenia, dopóki ktoś nie wznowi korutyny.
\end{description}

Jeżeli \texttt{initial\_suspend()} zwraca \texttt{suspend\_always},
korutyna jest \emph{lazy}: nie rozpoczyna wykonywania ciała przy samym
wywołaniu funkcji. Natomiast \texttt{suspend\_never} oznacza
\emph{eager start}: ciało może rozpocząć się natychmiast w wątku
wywołującym i pracować aż do pierwszego rzeczywistego zawieszenia.

\textbf{Ważne:} nie wolno wywoływać \texttt{resume()} na korutynie,
która jest zawieszona w swoim \texttt{final\_suspend}. Jest to
niezdefiniowane zachowanie. Uchwyt do takiej ramki może nadal służyć
do \texttt{destroy()}, jeżeli ramka nie została już zniszczona.

\section{Uchwyt \texttt{std::coroutine\_handle} i własność ramki}

\texttt{std::coroutine\_handle<Promise>} jest lekkim uchwytem
umożliwiającym m.in. wznowienie korutyny i dostęp do jej obietnicy.
Sam \texttt{coroutine\_handle} \emph{nie zarządza automatycznie
czasem życia ramki}. Skopiowanie uchwytu nie jest kopią korutyny,
a zniszczenie uchwytu nie wywołuje automatycznie \texttt{destroy()}.
Dlatego wygodna klasa będąca właścicielem ramki powinna stosować RAII,
uniemożliwiać kopiowanie oraz poprawnie obsługiwać przenoszenie.

\begin{lstlisting}[caption={Operacje na uchwycie; fragment ilustracyjny.}]
#include <coroutine>

using Handle = std::coroutine_handle<>;

// Zakladamy, ze h wskazuje na istniejaca, zawieszona korutyne.
// if (h && !h.done()) h.resume();
// h.destroy(); // dokladnie raz, przez wlasciciela ramki
// Po destroy() wszystkie kopie h sa wiszacymi uchwytami.
\end{lstlisting}

Uchwyt może zostać przekazany do innego wątku lub executora, ale
wówczas trzeba zagwarantować, że ramka nie zostanie zniszczona,
zanim planista faktycznie wznowi zawieszoną korutynę. Jest to jedna
z głównych trudności implementacyjnych asynchronicznego API C++20.

\section{Własny generator: \texttt{co\_yield} od podstaw}
\label{sec:cpp15-generator}

Zbudujemy teraz typ \texttt{Generator}, który produkuje liczby całkowite
kolejno, \emph{na żądanie}. Poniższy przykład jest samowystarczalny.
Prywatny uchwyt do ramki stanowi jedynego właściciela zasobu;
\texttt{next()} wznawia korutynę tylko wtedy, gdy nie zakończyła
wykonywania. Własną obsługą wyjątku uzupełniamy przykład o
poprawną propagację błędów.

\begin{lstlisting}[caption={Generator liczb całkowitych: co\_yield i RAII (C++20); kompletny program.}]
#include <cassert>
#include <coroutine>
#include <exception>
#include <limits>
#include <stdexcept>
#include <utility>

class Generator {
public:
    struct promise_type {
        int current = 0;
        std::exception_ptr error;

        Generator get_return_object() {
            return Generator{Handle::from_promise(*this)};
        }
        std::suspend_always initial_suspend() noexcept { return {}; }
        std::suspend_always final_suspend() noexcept { return {}; }
        std::suspend_always yield_value(int x) noexcept {
            current = x;
            return {};
        }
        void return_void() noexcept {}
        void unhandled_exception() noexcept {
            error = std::current_exception();
        }
    };

    using Handle = std::coroutine_handle<promise_type>;

    explicit Generator(Handle h) noexcept : handle_(h) {}
    Generator(const Generator&) = delete;
    Generator& operator=(const Generator&) = delete;

    Generator(Generator&& other) noexcept
        : handle_(std::exchange(other.handle_, {})),
          ready_(std::exchange(other.ready_, false)) {}

    Generator& operator=(Generator&& other) noexcept {
        if (this != &other) {
            if (handle_) handle_.destroy();
            handle_ = std::exchange(other.handle_, {});
            ready_ = std::exchange(other.ready_, false);
        }
        return *this;
    }

    ~Generator() {
        if (handle_) handle_.destroy();
    }

    bool next() {
        ready_ = false;
        if (!handle_ || handle_.done()) return false;
        handle_.resume();
        if (handle_.promise().error)
            std::rethrow_exception(handle_.promise().error);
        ready_ = !handle_.done();
        return ready_;
    }

    int value() const {
        if (!ready_)
            throw std::logic_error("No current value");
        return handle_.promise().current;
    }

private:
    Handle handle_{};
    bool ready_ = false;
};

Generator squares(int n) {
    for (int i = 1; i <= n; ++i) {
        if (i > std::numeric_limits<int>::max() / i)
            throw std::overflow_error("square outside int");
        co_yield i * i;
    }
}

int main() {
    auto gen = squares(4); // Cialo jeszcze sie nie wykonalo.
    int sum = 0;
    while (gen.next()) sum += gen.value();
    assert(sum == 1 + 4 + 9 + 16);
    assert(!gen.next());
}
\end{lstlisting}

\subsection{Co dzieje się podczas \texttt{co\_yield}?}

Instrukcja \texttt{co\_yield i * i} powoduje wywołanie
\texttt{promise.yield\_value(i * i)} i zastosowanie mechanizmu
\texttt{co\_await} do otrzymanego awaitera. W naszym przypadku
\texttt{yield\_value} zapisuje wartość w obietnicy i zwraca
\texttt{suspend\_always}. Funkcja \texttt{next()} otrzymuje wtedy
sterowanie z powrotem i udostępnia zapisany element. Następne
\texttt{next()} kontynuuje pętlę od poprzedniego miejsca zawieszenia.

Generator posiada tylko jeden bieżący wynik. Nie przechowuje
całego ciągu jak \texttt{std::vector<int>}. To użyteczne przy dużych
lub potencjalnie nieskończonych sekwencjach. \emph{Nie oznacza to}
jednak automatycznie mniejszego kosztu obliczeń --- dochodzą koszty
ramki i kolejnych wznowień.

\nota{Przedstawiony \texttt{Generator} jest celowo minimalny: nie
udostępnia interfejsu iteratorów \texttt{begin()/end()}, nie jest
bezpieczny do jednoczesnego wznawiania z kilku wątków i nie powinien
być stosowany jako gotowa biblioteka produkcyjna.}

\section{Protokół \texttt{co\_await}: \texttt{await\_ready},
\texttt{await\_suspend}, \texttt{await\_resume}}
\label{sec:cpp15-awaiter}

Operator \texttt{co\_await} nie jest synonimem oczekiwania na
\texttt{future}. Kompilator uzyskuje obiekt \emph{awaiter}, którego
interfejs ma trzy fundamentalne operacje:
\begin{description}
\item[\texttt{await\_ready()}] Zwraca wartość logiczną: jeśli
      \texttt{true}, zawieszenie jest pomijane.
\item[\texttt{await\_suspend(handle)}] Jeżeli \texttt{await\_ready()}
      było fałszywe, metoda podejmuje decyzję o kontynuacji.
      Dopuszczalne postacie zwracają \texttt{void}, \texttt{bool}
      albo inny \texttt{std::coroutine\_handle}.
\item[\texttt{await\_resume()}] Wyznacza wynik całego wyrażenia
      \texttt{co\_await}; może też zgłosić wyjątek.
\end{description}

Gdy \texttt{await\_suspend} zwraca \texttt{bool}, wartość
\texttt{false} oznacza kontynuację bez pozostawania w zawieszeniu,
a \texttt{true} --- pozostawienie korutyny w stanie zawieszonym.
Przy zwrocie \texttt{void} korutyna pozostaje zawieszona do późniejszego
wznowienia. Zwrócenie innego uchwytu powoduje \emph{symmetric transfer}:
sterowanie przechodzi do wskazanej korutyny.

Nie każdy typ bezpośrednio dostarcza trzy wymienione metody: język
obsługuje też przeciążenie \texttt{operator co\_await} oraz
\texttt{await\_transform} zdefiniowane w \texttt{promise\_type}.
Sposób uzyskania awaitera zależy od tych reguł.

\subsection{Czysto lokalne zawieszenie bez nowego wątku}

Poniższy typ liczy wznowienia. Wykorzystamy go z prostą korutyną,
której właściciel ręcznie wywołuje \texttt{resume()}.

\begin{lstlisting}[caption={co\_await z własnym awaiterem: żadnego nowego wątku (C++20); kompletny program.}]
#include <cassert>
#include <coroutine>
#include <exception>
#include <utility>

struct Stepper {
    struct promise_type {
        Stepper get_return_object() {
            return Stepper{handle::from_promise(*this)};
        }
        std::suspend_always initial_suspend() noexcept { return {}; }
        std::suspend_always final_suspend() noexcept { return {}; }
        void return_void() noexcept {}
        void unhandled_exception() { std::terminate(); }
    };

    using handle = std::coroutine_handle<promise_type>;
    handle h;
    explicit Stepper(handle x) : h(x) {}
    Stepper(const Stepper&) = delete;
    Stepper(Stepper&& other) noexcept
        : h(std::exchange(other.h, {})) {}
    ~Stepper() { if (h) h.destroy(); }
    bool resume() {
        if (!h || h.done()) return false;
        h.resume();
        return !h.done();
    }
};

struct Pause {
    int& count;
    bool await_ready() const noexcept { return false; }
    void await_suspend(std::coroutine_handle<>) const noexcept {
        ++count; // Zawiesz, az wlasciciel wywola resume().
    }
    void await_resume() const noexcept {}
};

Stepper example(int& count) {
    ++count;                 // wykona sie przy pierwszym resume
    co_await Pause{count};  // pierwsze zawieszenie
    ++count;                 // wykona sie przy drugim resume
}

int main() {
    int count = 0;
    auto step = example(count);
    assert(count == 0);
    assert(step.resume() && count == 2);
    assert(!step.resume() && count == 3);
}
\end{lstlisting}

Kod niczego nie oblicza równolegle. Wszystkie instrukcje wykonał
wątek \texttt{main}. Korutyna jedynie oddała sterowanie, zachowując
miejsce wznowienia.

\section{Własne \texttt{Task<T>}: leniwe obliczenie a asynchroniczność}
\label{sec:cpp15-task}

Zaprojektujemy najmniejszy rozsądny typ \texttt{Task<T>} zwracający
wynik i obsługujący wyjątek. Celowo nie implementujemy jeszcze
harmonogramu: pozwala to zrozumieć, że \emph{sam typ korutyny nie jest
schedulerem}. \texttt{get()} w tym przykładzie działa wyłącznie
wtedy, gdy ciało może zakończyć się po pojedynczym wznowieniu.

\begin{lstlisting}[caption={Lazy Task<T> bez schedulera (C++20); kompletny program.}]
#include <cassert>
#include <coroutine>
#include <exception>
#include <optional>
#include <stdexcept>
#include <utility>

template<class T>
class Task {
public:
    struct promise_type {
        std::optional<T> result;
        std::exception_ptr error;
        Task get_return_object() {
            return Task{Handle::from_promise(*this)};
        }
        std::suspend_always initial_suspend() noexcept { return {}; }
        std::suspend_always final_suspend() noexcept { return {}; }
        void return_value(T x) { result.emplace(std::move(x)); }
        void unhandled_exception() noexcept {
            error = std::current_exception();
        }
    };
    using Handle = std::coroutine_handle<promise_type>;

    explicit Task(Handle h) noexcept : h_(h) {}
    Task(const Task&) = delete;
    Task& operator=(const Task&) = delete;
    Task(Task&& x) noexcept
        : h_(std::exchange(x.h_, {})),
          taken_(std::exchange(x.taken_, false)) {}
    ~Task() { if (h_) h_.destroy(); }

    T get() {
        if (!h_) throw std::logic_error("Empty task");
        if (taken_) throw std::logic_error("Result already consumed");
        if (!h_.done()) h_.resume();
        if (!h_.done())
            throw std::logic_error("Task needs an external scheduler");
        taken_ = true;
        if (h_.promise().error)
            std::rethrow_exception(h_.promise().error);
        return std::move(*h_.promise().result);
    }

private:
    Handle h_{};
    bool taken_ = false;
};

Task<int> calculate(int x) {
    co_return x * x;
}

int main() {
    auto task = calculate(7); // lazy, obliczenie jeszcze nie ruszylo
    assert(task.get() == 49);
}
\end{lstlisting}

Ten \texttt{Task<T>} jest pedagogiczny: nie ma przeciążenia
\texttt{operator co\_await}, nie wspiera wielokrotnego pobierania
wyniku (drugie \texttt{get()} zgłasza błąd) i nie potrafi samodzielnie
czekać na zdarzenia zewnętrzne.
W szczególności \texttt{get()} nie może po prostu wielokrotnie
wywoływać \texttt{resume()} na rzeczywistym zadaniu I/O: wznowienie
przed ukończeniem operacji mogłoby naruszyć protokół awaitera.
Produkcyjne \texttt{Task<T>} wymaga m.in. protokołu kontynuacji,
jasnej semantyki własności i poprawnej synchronizacji z executorem.

\section{Korutyna a \texttt{std::async}: fundamentalne porównanie}
\label{sec:cpp15-async-comparison}

\texttt{std::async} (C++11, nagłówek \texttt{<future>}) jest gotową
funkcją uruchamiającą zadanie i zwracającą \texttt{std::future<T>}.
Nie wymaga od użytkownika implementowania \texttt{promise\_type}.
Przy polityce \texttt{std::launch::async} funkcja pracuje w odrębnym
wątku; polityka \texttt{std::launch::deferred} odracza jej uruchomienie
do pierwszego nieczasowego oczekiwania na wynik.

Z kolei korutyna C++20 definiuje sposób zapisywania \emph{kontynuacji}.
Nawet jeśli wynik ma charakter asynchroniczny, \emph{faktyczne
uruchamianie pracy} zależy od awaitera, event loop, schedulera lub
innego mechanizmu zewnętrznego. Nie istnieje standardowa ogólna
reguła ,,jedna korutyna = jeden wątek''.


\paragraph{Czy można napisać \texttt{co\_await future}?}
Nie wprost: w C++20 standardowy \texttt{std::future<T>} nie udostępnia
protokołu awaitera wymaganego przez \texttt{co\_await}.
\texttt{future.get()} jest operacją potencjalnie blokującą
bieżący wątek. Istnieją biblioteki dostarczające awaitery dla
przyszłych wyników, ale sam adapter nie zamienia blokującej
operacji w nieblokującą. Naiwne uruchamianie osobnego wątku
tylko po to, by wykonał \texttt{future.wait()} i wznowił korutynę,
może zniweczyć zalety oszczędnego używania wątków.
Dla rzeczywiście nieblokującego oczekiwania na I/O potrzeba
integracji z właściwym mechanizmem asynchronicznym.

\begin{center}
\small
\begin{tabularx}{\textwidth}{@{}lXX@{}}
\toprule
\textbf{Cecha} & \textbf{\texttt{std::async}} & \textbf{Korutyny C++20} \\
\midrule
Co opisuje? & Zadanie ze stanem wyniku & Wykonanie z punktami zawieszenia \\
Tworzy nowy wątek? & \texttt{launch::async}: osobny wątek & Samo \texttt{co\_await}: nie \\
Kiedy zaczyna pracę? & Zależnie od polityki & Zależnie od \texttt{initial\_suspend} i schedulera \\
Odbieranie wyniku & \texttt{future::get()} (może blokować) & Zależy od typu \texttt{Task} i awaitera \\
Wyjątki & Przechowywane w stanie \texttt{future} & Obsługa przez \texttt{unhandled\_exception} i typ wyniku \\
Anulowanie & Brak automatycznego anulowania przez \texttt{future} & Możliwe kooperacyjnie; wymaga protokołu \\
Pamięć/zasoby & Stan asynchroniczny, przy \texttt{async} również wątek & Ramka korutyny, ewentualnie executor \\
Gotowe narzędzie w C++20? & Tak & Składnia tak; ogólny \texttt{Task<T>} nie \\
\bottomrule
\end{tabularx}
\end{center}

\subsection{Ten sam rachunek zapisany na dwa sposoby}

Porównanie dwóch krótkich funkcji pokazuje różnicę semantyczną.

\begin{lstlisting}[caption={std::async uruchamia zadanie w odrębnym wątku (C++11); kompletny program.}]
#include <cassert>
#include <future>

int square(int x) { return x * x; }

int main() {
    auto f = std::async(std::launch::async, square, 7);
    assert(f.get() == 49);
}
\end{lstlisting}

\begin{lstlisting}[caption={Korutyna bez awaitera zewnętrznego nie tworzy wątku; fragment.}]
// Przyjmujemy implementacje Task<T> z poprzedniej sekcji.
Task<int> square_coroutine(int x) {
    co_return x * x;
}

// Task<int> t = square_coroutine(7);
// int result = t.get(); // 49, w watku wywolujacym get()
\end{lstlisting}

W obu przypadkach można otrzymać tę samą wartość, ale nie jest to
to samo zachowanie wykonawcze. \texttt{std::async} z \texttt{launch::async}
zarządza wykonaniem w innym wątku, a pokazany \texttt{Task}
jedynie odkłada wywołanie do momentu ręcznego wznowienia.

\subsection{Pułapka destruktora \texttt{future}}

Ważna różnica dotyczy czasu życia. Gdy ostatni uchwyt do stanu
asynchronicznego utworzonego przez \texttt{std::async} z polityką
\texttt{launch::async} jest niszczony i praca wciąż trwa, destruktor
czeka na zakończenie powiązanego wątku.
Pierwsze z poniższych wywołań zakończy się, zanim rozpocznie się drugie:

\begin{lstlisting}[caption={Zignorowany future wymusza oczekiwanie; fragment.}]
// Funkcje slow_a i slow_b wykonuja dlugie obliczenia.
// std::async(std::launch::async, slow_a);
// std::async(std::launch::async, slow_b);
\end{lstlisting}

Jeżeli chcemy umożliwić nakładanie się zadań, trzeba przechować ich
uchwyty przed oczekiwaniem:
\begin{lstlisting}[caption={Dwa przyszłe wyniki powinny żyć jednocześnie; fragment.}]
// auto fa = std::async(std::launch::async, slow_a);
// auto fb = std::async(std::launch::async, slow_b);
// auto a = fa.get();
// auto b = fb.get();
\end{lstlisting}

Analogicznie w systemie korutyn należy znać politykę niszczenia
\texttt{Task}: usunięcie właściciela zawieszonej korutyny bez
koordynacji z executorem może zostawić wiszący uchwyt. Różne
biblioteki oferują odmienne kontrakty anulowania i oczekiwania.

\subsection{Asynchroniczne I/O a zadania CPU-bound}

Przy \emph{operacji I/O} większość czasu może upływać na oczekiwaniu
na dysk, sieć lub timer. Duża liczba korutyn może opisywać wiele
rozpoczętych operacji, podczas gdy mała pula wątków obsługuje
zdarzenia gotowości i wznowienia. Korzyść nie bierze się z tego,
że \texttt{co\_await} ,,jest szybsze'', lecz z możliwości
nieblokowania wątków podczas oczekiwania.

Przy \emph{operacjach CPU-bound} (obliczenia numeryczne, kompresja,
przetwarzanie obrazów) samo zawieszenie korutyny nie zwiększa
liczby rdzeni wykonujących kod. Tu potrzebne jest rzeczywiste
rozłożenie pracy na wiele wątków (np. \texttt{std::async}, pula
wątków lub OpenMP). Korutyna może być wygodnym \emph{interfejsem}
do takiego executora, ale nie zastępuje go.

\section{Rzeczywista asynchroniczność: korutyna i własny executor}
\label{sec:cpp15-executor}

Żeby zestawienie z \texttt{std::async} nie pozostało czysto
teoretyczne, utworzymy prosty executor z jednym wątkiem roboczym.
\texttt{co\_await executor.schedule()} odłoży kontynuację do
kolejki. Wznowienie wykona \emph{wątek executora}, nie wywołujący.
Wartością zwracaną będzie nasza klasa \texttt{FutureTask<T>},
wykorzystująca wewnętrznie \texttt{std::promise<T>} i
\texttt{std::future<T>}.

W przeciwieństwie do poprzedniego przykładu ramka ma strategię
\emph{samozniszczenia} przez \texttt{final\_suspend = suspend\_never}.
Nie przechowujemy już zewnętrznego uchwytu właścicielskiego:
po zakończeniu korutyna sama domyka swój cykl życia, a klient posiada
wyłącznie \texttt{future}. W tej uproszczonej konstrukcji executor
musi żyć aż do zakończenia zadań, a pojedyncza korutyna jest
wznawiana w kolejce co najwyżej raz.

\begin{lstlisting}[caption={Korutyna przechodząca do innego wątku poprzez executor (C++20); kompletny program.}]
#include <cassert>
#include <condition_variable>
#include <coroutine>
#include <deque>
#include <exception>
#include <future>
#include <mutex>
#include <stdexcept>
#include <thread>
#include <utility>

class Executor {
    std::mutex m_;
    std::condition_variable cv_;
    std::deque<std::coroutine_handle<>> queue_;
    bool closed_ = false;
    std::thread worker_;

    void run() {
        for (;;) {
            std::coroutine_handle<> h;
            {
                std::unique_lock lock(m_);
                cv_.wait(lock, [&] { return closed_ || !queue_.empty(); });
                if (queue_.empty() && closed_) return;
                h = queue_.front();
                queue_.pop_front();
            }
            h.resume(); // H moze sie tu samozniszczyc.
            // Nie korzystamy juz z h po resume().
        }
    }

public:
    Executor() : worker_([this] { run(); }) {}
    Executor(const Executor&) = delete;
    Executor& operator=(const Executor&) = delete;

    ~Executor() {
        {
            std::lock_guard lock(m_);
            closed_ = true;
        }
        cv_.notify_one();
        worker_.join();
    }

    void post(std::coroutine_handle<> h) {
        {
            std::lock_guard lock(m_);
            if (closed_) throw std::runtime_error("Executor closed");
            queue_.push_back(h);
        }
        cv_.notify_one();
    }

    struct Schedule {
        Executor& ex;
        bool await_ready() const noexcept { return false; }
        void await_suspend(std::coroutine_handle<> h) const {
            ex.post(h); // Po przekazaniu h nie uzywamy juz awaitera.
        }
        void await_resume() const noexcept {}
    };

    Schedule schedule() noexcept { return Schedule{*this}; }
};

template<class T>
class FutureTask {
    std::future<T> future_;
public:
    explicit FutureTask(std::future<T> f) : future_(std::move(f)) {}
    FutureTask(FutureTask&&) = default;
    FutureTask(const FutureTask&) = delete;

    T get() { return future_.get(); }

    struct promise_type {
        std::promise<T> result;
        FutureTask get_return_object() {
            return FutureTask{result.get_future()};
        }
        std::suspend_never initial_suspend() noexcept { return {}; }
        std::suspend_never final_suspend() noexcept { return {}; }
        void return_value(T x) { result.set_value(std::move(x)); }
        void unhandled_exception() noexcept {
            try { result.set_exception(std::current_exception()); }
            catch (...) { std::terminate(); }
        }
    };
};

FutureTask<int> square_on(Executor& ex, int x) {
    co_await ex.schedule(); // Tu oddajemy sterowanie.
    co_return x * x;        // Wykona sie w watku roboczym.
}

int main() {
    Executor ex;
    auto a = square_on(ex, 6);
    auto b = square_on(ex, 7);
    assert(a.get() + b.get() == 85);
    // Przed zniszczeniem ex oba zadania sie zakonczyly.
}
\end{lstlisting}

\paragraph{Co to porównanie rzeczywiście dowodzi?}
Kod z \texttt{square\_on} jest asynchroniczny względem wywołującego,
ponieważ po \texttt{co\_await} kontynuację przejmuje inny wątek.
Nie oznacza to jeszcze \emph{równoległego wykonywania} obu obliczeń:
executor ma tylko jednego pracownika, więc zadania są przez niego
wykonywane kolejno. Do równoległego przetwarzania dwóch obliczeń
należałoby dodać drugiego pracownika albo wykorzystać istniejącą
pulę wątków.

Pokazany executor jest \emph{szkicem edukacyjnym}. Nie obsługuje
ogólnego anulowania ani ponownego planowania tej samej korutyny,
nie rozwiązuje wszystkich problemów czasu życia i nie stanowi
gotowej biblioteki. W szczególności klient nie ma uchwytu pozwalającego
anulować i bezpiecznie zniszczyć zawieszone zadanie. W systemie
produkcyjnym własność korutyny, zamknięcie executora i czas życia
powiązanych zasobów muszą stanowić jeden, spójny kontrakt.

\section{Obsługa wyjątków, anulowanie i czas życia danych}
\label{sec:cpp15-errors}

\subsection{Wyjątki a \texttt{unhandled\_exception}}

Nieobsłużony wyjątek z ciała korutyny nie może zostać pozostawiony
bez decyzji. Kompilator wywołuje \texttt{promise.unhandled\_exception()}.
W generatorze zapisaliśmy \texttt{std::exception\_ptr}; przy kolejnym
\texttt{next()} błąd zostanie zgłoszony w kodzie klienta. W naszym
\texttt{FutureTask<T>} wyjątek trafia do stanu \texttt{std::future}
i ujawnia się podczas \texttt{get()}. Jest to analogiczne do
propagacji błędów w \texttt{std::async}, ale sam mechanizm wymagał
jawnej implementacji w \texttt{promise\_type}.

Należy też uwzględnić wyjątki zgłaszane przez operacje awaitera
oraz błędy alokacji ramki. Awaria podczas tworzenia ramki może
wystąpić jeszcze przed wykonaniem ciała korutyny. Typ obietnicy
może opcjonalnie udostępniać \texttt{get\_return\_object\_on\_allocation\_failure()}
razem z odpowiednią, niezgłaszającą wyjątku alokacją, ale nie jest to
mechanizm automatycznie dostępny każdemu typowi korutyny.

\subsection{Referencje przechowywane w ramce}

Argumenty referencyjne nie wydłużają czasu życia obiektów.
Jeżeli korutyna przechowuje referencję do zmiennej lokalnej, która
zniknie podczas zawieszenia, jej późniejsze wznowienie może prowadzić
do użycia nieistniejącego już obiektu.

\begin{lstlisting}[caption={Niebezpieczny czas życia referencji w korutynie; fragment koncepcyjny.}]
// Zakladamy lazy Task<int> z poprzedniej sekcji.
Task<int> read_later(const int& ref) {
    co_return ref;
}

// NIEPOPRAWNE:
// auto task = read_later(123); // 123 znika na koncu tej instrukcji
// int value = task.get();     // odczyt przez wiszaca referencje
\end{lstlisting}

Z reguły bezpieczniej przechowywać wewnątrz ramki wartości
przeznaczone do późniejszego użycia (np. obiekt przekazany przez
wartość), a nie referencje do tymczasowych obiektów.
Trzeba też zwracać uwagę na czas życia domknięć lambda:
korutyna będąca lambdą i wykorzystująca przechwycenia może
po zawieszeniu odwoływać się do już zniszczonego obiektu domknięcia.
Sam fakt przeniesienia wynikowej korutyny nie musi temu zapobiegać.

\subsection{Anulowanie jest kooperacyjne}

\texttt{std::stop\_token} (C++20) umożliwia \emph{żądanie}, a nie
wymuszenie zakończenia zadania. Korutyna może sprawdzać token
przed rozpoczęciem kolejnych faz albo korzystać z awaitera,
który potrafi odpowiednio reagować na anulowanie. Nie wynika
z tego jednak, że każde \texttt{co\_await} automatycznie obsługuje
\texttt{stop\_token}.

\begin{lstlisting}[caption={Kooperacyjne anulowanie korutyny: sprawdzenie żądania w punktach kontrolnych; fragment.}]
// Zakladamy Task<int> z poprzedniej sekcji.
#include <stop_token>

Task<int> calculate(std::stop_token token) {
    if (token.stop_requested()) co_return -1;
    // Pierwszy krok obliczen...
    if (token.stop_requested()) co_return -1;
    // Kolejny krok obliczen...
    co_return 42;
}
\end{lstlisting}

W powyższym przykładzie token jest przekazywany przez wartość.
Sam \texttt{stop\_token} nie steruje czasem życia ramki i nie
wznawia korutyny, która czeka na zdarzenie. Do pełnego anulowania
operacji I/O potrzeba jeszcze protokołu z usługą, która wykonuje
rzeczywistą operację, oraz reguł określających \emph{kto} zniszczy
korutynę i \emph{kiedy} wolno to zrobić.

\section{Pułapki synchronizacji i planowania}

\subsection{Równoczesne wznowienie tej samej korutyny}

Nie można traktować \texttt{coroutine\_handle::resume()} jak
operacji atomowej nadającej się do równoczesnego wywoływania przez
wiele wątków. Trzeba zapewnić, że dana zawieszona korutyna ma
jednego logicznego właściciela uprawnienia do wznowienia.
Ponowne wznowienie korutyny już aktywnej lub zakończonej może
prowadzić do niezdefiniowanego zachowania.

\subsection{Utracone powiadomienie i wyścig publikacji uchwytu}

Awaiter może przekazać uchwyt do innego wątku wewnątrz
\texttt{await\_suspend}. Wówczas druga strona może wznowić korutynę
\emph{zanim} pierwotne wywołanie \texttt{await\_suspend} powróci.
Dlatego po opublikowaniu uchwytu kod awaitera nie powinien korzystać
z obiektów, których czas życia mógł się już zakończyć.

W bardziej złożonych awaiterach trzeba również uniknąć zgubienia
zdarzenia kończącego operację między sprawdzeniem gotowości a
zarejestrowaniem oczekiwania. Poprawne rozwiązania stosują np.
mutex, operacje atomowe albo protokół jednorazowego przejęcia
kontynuacji. Samo \texttt{await\_ready()} nie zastępuje synchronizacji.

\subsection{Brak własności \texttt{coroutine\_handle}}

Gdy planista ma w kolejce jedynie \texttt{coroutine\_handle<>},
nie posiada informacji, kto może zniszczyć ramkę. Usunięcie
obiektu \texttt{Task} podczas oczekiwania może prowadzić do
późniejszego wznowienia już zwolnionej pamięci.
Najbezpieczniejszy projekt definiuje własność od początku:
\emph{Task posiada ramkę}, \emph{executor posiada ramkę} albo
\emph{ramka posiada samą siebie}, z dokładnym protokołem przekazania
własności i zakończenia.

\section{Koszty, optymalizacje i pomiary}
\label{sec:cpp15-performance}

Wydajność należy rozpatrywać w całym systemie, a nie tylko przez
porównanie czasu jednego \texttt{co\_await} i jednego \texttt{std::async}.
\begin{itemize}
\item \textbf{Utworzenie wątku.} \texttt{std::async(launch::async, ...)}
      uruchamia pracę w odrębnym wątku; bardzo duża liczba krótkich
      zadań może generować istotny narzut zarządzania wątkami.
\item \textbf{Ramka.} Korutyna może wymagać alokacji pamięci,
      zachowywania stanu i wznawiania. Kompilator może czasem
      wyeliminować alokację ramki, ale nie jest to gwarancja.
\item \textbf{Executor.} Kolejka, mutexy, przełączanie między wątkami
      oraz współdzielenie pamięci również kosztują.
\item \textbf{Lokalność pamięci.} Duża liczba korutyn z rozproszonymi
      ramkami może prowadzić do dodatkowych nietrafień cache.
\item \textbf{Skalowalność.} Dla I/O ważniejsza może być liczba
      oczekujących operacji niż przepustowość pojedynczego zadania.
      Dla CPU-bound kluczowa jest liczba rdzeni, wielkość porcji pracy
      i ograniczenie dostępu do współdzielonych zasobów.
\end{itemize}

Rzetelny benchmark powinien mierzyć co najmniej: całkowitą
przepustowość, opóźnienia (również wartości odstające), liczbę wątków,
liczbę alokacji oraz maksymalne zużycie pamięci. Porównywanie
\texttt{std::async} z \emph{nieuruchamiającą nowych wątków} korutyną
wykonującą natychmiast prostą arytmetykę nie jest miarodajnym
testem wydajności asynchronicznej.

\subsection{Praktyczna decyzja: kiedy używać którego narzędzia?}

\begin{center}
\small
\begin{tabularx}{\textwidth}{@{}lX@{}}
\toprule
\textbf{Problem} & \textbf{Rozsądny pierwszy wybór} \\
\midrule
Jedno niezależne zadanie CPU z wynikiem & \texttt{std::async(std::launch::async, ...)} \\
Wiele krótkich zadań CPU & Pula wątków, executor lub OpenMP \\
Wiele połączeń I/O z długim oczekiwaniem & Korutyny + asynchroniczny runtime I/O \\
Leniwe generowanie sekwencji & Generator korutynowy \\
Wspólny licznik lub krótka sekcja krytyczna & \texttt{atomic} albo mutex, nie korutyna \\
Anulowalne zadanie działające długo & Przemyślany protokół \texttt{stop\_token} \\
Interfejs biblioteki o złożonych zależnościach & \texttt{Task<T>} nad zdefiniowanym executorem \\
\bottomrule
\end{tabularx}
\end{center}

\section{Związek z OpenMP, \texttt{thread} i \texttt{future}}

Rozszerzony rozdział o współbieżności omawiał \texttt{std::thread},
\texttt{std::jthread}, atomiki, blokady, \texttt{std::async} oraz OpenMP.
Korutyny nie unieważniają tych narzędzi; stanowią kolejną warstwę.
\begin{itemize}
\item \texttt{std::thread} opisuje konkretny wątek wykonania,
      a \texttt{std::jthread} dodatkowo pomaga zarządzać jego
      czasem życia i kooperacyjnym zatrzymywaniem.
\item \texttt{std::async} tworzy stan przyszłego wyniku i
      uruchamia wywołanie zgodnie z przyjętą polityką.
\item OpenMP rozdziela pętle i zadania między wątki obsługiwane
      przez środowisko wykonawcze OpenMP; szczególnie przydaje się
      w obliczeniach CPU-bound i kodzie numerycznym.
\item Korutyny pozwalają wygodnie opisać wieloetapową logikę
      oczekiwania i kontynuacji. Ich wydajność i stopień równoległości
      zależą od dostarczonego środowiska wykonawczego.
\end{itemize}

Możliwe jest użycie \emph{korutyn nad pulą wątków}: \texttt{co\_await}
przekazuje kontynuację do kolejki, a pracownik puli później ją wznawia.
Nie należy jednak od razu łączyć wszystkich mechanizmów.
Najpierw określamy, czy program czeka głównie na I/O, czy wykonuje
obliczenia, a następnie wybieramy model wykonania i własności zasobów.

\section{Najczęstsze błędy projektowe}

\begin{enumerate}
\item Zakładanie, że \texttt{co\_await} uruchamia nowy wątek.
\item Próba zwrócenia bezpośrednio \texttt{std::future<T>}
      z własnej korutyny bez odpowiedniego \texttt{promise\_type}.
\item Dwukrotne wywołanie \texttt{destroy()} dla tej samej ramki.
\item Wywołanie \texttt{resume()} po osiągnięciu
      \texttt{final\_suspend}.
\item Pozostawienie referencji lub przechwycenia lambdy do obiektu,
      który znika podczas zawieszenia.
\item Zniszczenie \texttt{Task} mimo tego, że executor posiada
      w kolejce uchwyt do jego ramki.
\item Mylenie automatycznej propagacji wyjątków w \texttt{std::async}
      z obowiązkami implementatora \texttt{promise\_type}.
\item Traktowanie \texttt{std::stop\_token} jako mechanizmu
      przymusowego przerwania dowolnego \texttt{co\_await}.
\item Równoczesne wznawianie tej samej korutyny z kilku wątków.
\item Uznanie, że zawieszenie i wznowienie zawsze kosztuje mniej
      niż zwykłe wywołanie funkcji.
\item Uruchamianie \texttt{std::async} dla tysięcy mikro-zadań
      bez rozważenia puli wątków.
\item Zwracanie z \texttt{await\_suspend} po opublikowaniu uchwytu
      i dalsze używanie potencjalnie zniszczonego awaitera.
\end{enumerate}

\section*{Podsumowanie rozdziału}
\addcontentsline{toc}{section}{Podsumowanie rozdziału}

Korutyny C++20 są mechanizmem \emph{zawieszenia i kontynuacji},
a nie samodzielnym środowiskiem asynchronicznym. Generator używa
\texttt{co\_yield}, aby leniwie udostępniać wartości. Typ zadania
używa \texttt{co\_return}, obiektu \texttt{promise\_type} i awaiterów,
aby opisać wykonanie z wynikiem i wyjątkami. \texttt{co\_await}
przekazuje kontrolę zgodnie z protokołem awaitera, a dopiero executor
lub inne środowisko może wznowić kod w osobnym wątku.

\texttt{std::async} jest prostym narzędziem do uruchomienia funkcji
i odebrania przyszłego wyniku. Korutyny są lepszym budulcem dla
złożonych przepływów wieloetapowych, szczególnie I/O, ale wymagają
większej odpowiedzialności za pamięć, synchronizację, anulowanie
i czas życia ramki. W obu przypadkach należy rozdzielać trzy pytania:
\emph{kto wykonuje kod, kto posiada stan i kto decyduje o zakończeniu}.

\section*{Literatura i dokumentacja}
\addcontentsline{toc}{section}{Literatura i dokumentacja}
\begin{itemize}
\item \url{https://en.cppreference.com/w/cpp/language/coroutines} --- model korutyn C++20.
\item \url{https://en.cppreference.com/w/cpp/coroutine/coroutine_handle} --- uchwyty i zarządzanie ramką.
\item \url{https://en.cppreference.com/w/cpp/coroutine/suspend_always} --- podstawowe awaitery.
\item \url{https://en.cppreference.com/w/cpp/thread/async} --- \texttt{std::async} i polityki uruchamiania.
\item \url{https://en.cppreference.com/w/cpp/thread/future} --- \texttt{std::future}.
\item \url{https://en.cppreference.com/w/cpp/coroutine/generator} --- \texttt{std::generator} (C++23).
\item Lewis Baker, \emph{C++ Coroutines: Understanding the Compiler Transformations}
      --- seria publikacji o modelu wykonania korutyn.
\item David Mazières, \emph{My Tutorial and Take on C++20 Coroutines}, 2021.
\end{itemize}

% Listing oznaczony jako „kompletny program” mozna skompilowac oddzielnie.

\chapter{Korutyny, zadania asynchroniczne i \texttt{std::async}}
\label{chap:cpp15-coroutines}

\section{Po co C++ potrzebuje korutyn?}

Zwykła funkcja ma zasadniczo jeden punkt wejścia. Może wywoływać inne
funkcje, ale po powrocie oddaje sterowanie wywołującemu i nie posiada
mechanizmu późniejszego wznowienia dokładnie w miejscu, w którym
zakończyła wykonywanie. \emph{Korutyna} (ang. \emph{coroutine}) jest
funkcją, której wykonanie można \emph{zawiesić}, a następnie
\emph{wznowić}, zachowując wybrane zmienne lokalne i miejsce kontynuacji.

Korutyny C++20 wspierają między innymi:
\begin{itemize}
\item \emph{generatory} --- wytwarzanie kolejnych wartości na żądanie,
      bez budowania całej kolekcji;
\item \emph{zadania asynchroniczne} --- opis operacji, których kolejne
      etapy uruchamiane są po zakończeniu wcześniejszych;
\item \emph{kooperacyjne planowanie} --- oddawanie sterowania
      harmonogramowi zadań;
\item \emph{maszyny stanów} --- zachowywanie stanu pomiędzy wznowieniami
      bez ręcznego kodowania każdej możliwej kontynuacji.
\end{itemize}

Najważniejsza zasada brzmi:
\begin{center}
\emph{Korutyna nie jest wątkiem, a \texttt{co\_await} nie oznacza
automatycznie ,,uruchom w nowym wątku''.}
\end{center}

Korutyna określa \emph{jak zatrzymać i kontynuować obliczenie};
\texttt{std::async}, \texttt{std::thread} i pula wątków dotyczą również
\emph{gdzie i kiedy} wykonywać pracę. W szczególności kod korutyny
może działać od początku do końca w jednym wątku, nawet jeśli
wielokrotnie się zawiesza. Odwrotnie, korutyna może być wznawiana
przez różne wątki, jeżeli tak zdecyduje dostarczony mechanizm planowania.

\subsection{Trzy słowa kluczowe}

\begin{center}\small
\begin{tabularx}{\textwidth}{@{}p{0.32\textwidth}X@{}}
\toprule
Konstrukcja & Działanie w korutynie \\
\midrule
\code{co_await expr} & Przekazuje sterowanie obiektowi oczekującemu; może zawiesić korutynę bez blokowania wątku. \\
\code{co_yield value} & Udostępnia kolejną wartość i zwykle zawiesza generator do wznowienia. \\
\code{co_return value} & Kończy korutynę i przekazuje wynik zgodnie z protokołem jej typu zwracanego. \\
\code{await_ready()} & Określa, czy wynik jest już gotowy i można pominąć zawieszenie. \\
\code{await_suspend(h)} & Obsługuje zawieszenie i uchwyt \code{h}; może zaplanować późniejsze wznowienie. \\
\code{await_resume()} & Zwraca rezultat oczekiwania albo zgłasza zapisany wyjątek po wznowieniu. \\
\bottomrule
\end{tabularx}
\end{center}

C++20 wprowadził następujące konstrukcje składniowe:
\begin{description}
\item[\texttt{co\_await expr}] Zastosuj protokół oczekiwania do wyrażenia;
      zależnie od niego możesz zawiesić korutynę i później ją wznowić.
\item[\texttt{co\_yield expr}] Udostępnij kolejną wartość generatora;
      składniowo jest to szczególna postać oczekiwania, wykorzystująca
      metodę \texttt{yield\_value} obiektu \texttt{promise\_type}.
\item[\texttt{co\_return expr}] Zakończ korutynę i przekaż wynik
      przez odpowiednie \texttt{return\_value} albo \texttt{return\_void}.
\end{description}

Funkcja zawierająca jeden z tych operatorów (w odpowiednim kontekście)
rozpoznawana jest jako korutyna. Typ jej zwracanej wartości musi
udostępniać właściwy \emph{typ obietnicy} \texttt{promise\_type}
(lub odpowiednią specjalizację \texttt{std::coroutine\_traits}).
To wymóg mechanizmu korutyn, a nie znanego już \texttt{std::promise}.

\begin{lstlisting}[caption={Niepoprawna intuicja: std::future nie jest automatycznie typem korutyny (C++20); fragment celowo błędny.}]
#include <future>

// Standardowe std::future<int> nie zapewnia promise_type:
// std::future<int> calculate() { co_return 42; } // BLAD!
\end{lstlisting}

W C++20 biblioteka standardowa dostarcza \emph{mechanizm} korutyn
(nagłówek \texttt{<coroutine>}), ale nie dostarcza gotowego, ogólnego
\texttt{Task<T>} ani standardowego executora. Takie typy definiuje
biblioteka lub programista. \texttt{std::generator} jest rozwiązaniem
późniejszym, wprowadzonym dopiero w C++23, a jego wsparcie zależy
od użytej biblioteki standardowej.

\section{Różnica między zawieszeniem a zablokowaniem wątku}
\label{sec:cpp15-suspend-block}

Rozważmy dwa sposoby wstrzymania działania programu:
\begin{itemize}
\item \texttt{std::this\_thread::sleep\_for(100ms)} \emph{blokuje
      bieżący wątek} na pewien czas; sam wątek pozostaje zajętym
      zasobem systemowym, choć może nie zużywać wówczas czasu CPU;
\item dobrze zaprojektowane \texttt{co\_await} \emph{zawiesza
      korutynę}, oddając wątek do wykonywania innych zadań;
      wznowienie może nastąpić po nadejściu wyniku operacji I/O,
      komunikatu lub upływie czasu.
\end{itemize}

Nie każde \texttt{co\_await} daje zawieszenie: awaiter może zgłosić,
że wynik jest już gotowy. Również samo użycie \texttt{co\_await}
nie gwarantuje nieblokującego I/O --- zależy to od jego implementacji.

\begin{center}
\small
\begin{tabularx}{\textwidth}{@{}lXX@{}}
\toprule
\textbf{Właściwość} & \textbf{Blokada wątku} & \textbf{Zawieszenie korutyny} \\
\midrule
Zachowanie stanu & Na stosie wywołań wątku & Zwykle w ramce korutyny \\
Wątek potrzebny podczas oczekiwania & Tak, istnieje i czeka & Nie jest konieczny \emph{dla samego oczekiwania} \\
Wznowienie & Planista systemu/odblokowanie & Mechanizm awaitera i właściciel kontynuacji \\
Koszt & M.in. wątek i jego stos & M.in. ramka i obsługa zawieszenia \\
\bottomrule
\end{tabularx}
\end{center}

\section{Ramka korutyny i jej cykl życia}
\label{sec:cpp15-frame}

Kompilator przekształca korutynę w obiekt reprezentujący jej stan.
Ten stan nazywamy \emph{ramką korutyny} (ang. \emph{coroutine frame}).
Może ona obejmować obiekt \texttt{promise\_type}, parametry i stan
potrzebny do kontynuowania obliczenia po zawieszeniu. To, które
zmienne lokalne trzeba zachować, wynika z przebiegu programu.

W ogólnym modelu wywołanie korutyny przebiega w następujących etapach:
\begin{enumerate}
\item Alokowana jest ramka, chyba że implementacja legalnie
      wyeliminuje tę alokację (np. przez zagnieżdżenie ramki w innym obiekcie).
\item Kopiowane lub przenoszone są parametry odpowiednio do ich rodzajów;
      \emph{parametry referencyjne pozostają referencjami}.
\item Konstruowany jest obiekt obietnicy \texttt{promise\_type}.
\item Wywoływana jest metoda \texttt{get\_return\_object()}.
\item Wykonywany jest \texttt{co\_await promise.initial\_suspend()}.
\item Uruchamiane jest ciało korutyny (od razu lub po wznowieniu).
\item \texttt{co\_return} uruchamia \texttt{return\_value} lub
      \texttt{return\_void}; nieobsłużony wyjątek trafia do
      \texttt{unhandled\_exception()}.
\item Niszczone są obiekty lokalne, a korutyna wykonuje końcowy
      \texttt{co\_await} na wyniku \texttt{promise.final\_suspend()}.
\item Ramka zostaje ostatecznie zniszczona we właściwym momencie.
\end{enumerate}

Kolejność i czas \emph{zwolnienia ramki} zależą od przyjętej polityki
własności. Dla \texttt{final\_suspend = suspend\_always} ramkę zwykle
niszczy właściciel uchwytu, nawet po zakończeniu ciała korutyny.
Przy \texttt{final\_suspend = suspend\_never} mechanizm korutyny
może zakończyć się bez pozostawiania ramki zawieszonej do późniejszego
zniszczenia przez właściciela. Tych strategii nie wolno mieszać.

\subsection{Początek i koniec: \texttt{initial\_suspend} i \texttt{final\_suspend}}

Dwa najczęściej używane typy z nagłówka \texttt{<coroutine>} to:
\begin{description}
\item[\texttt{std::suspend\_never}] Oczekiwanie jest natychmiast
      gotowe i nie powoduje zawieszenia.
\item[\texttt{std::suspend\_always}] Oczekiwanie zawsze prowadzi
      do zawieszenia, dopóki ktoś nie wznowi korutyny.
\end{description}

Jeżeli \texttt{initial\_suspend()} zwraca \texttt{suspend\_always},
korutyna jest \emph{lazy}: nie rozpoczyna wykonywania ciała przy samym
wywołaniu funkcji. Natomiast \texttt{suspend\_never} oznacza
\emph{eager start}: ciało może rozpocząć się natychmiast w wątku
wywołującym i pracować aż do pierwszego rzeczywistego zawieszenia.

\textbf{Ważne:} nie wolno wywoływać \texttt{resume()} na korutynie,
która jest zawieszona w swoim \texttt{final\_suspend}. Jest to
niezdefiniowane zachowanie. Uchwyt do takiej ramki może nadal służyć
do \texttt{destroy()}, jeżeli ramka nie została już zniszczona.

\section{Uchwyt \texttt{std::coroutine\_handle} i własność ramki}

\texttt{std::coroutine\_handle<Promise>} jest lekkim uchwytem
umożliwiającym m.in. wznowienie korutyny i dostęp do jej obietnicy.
Sam \texttt{coroutine\_handle} \emph{nie zarządza automatycznie
czasem życia ramki}. Skopiowanie uchwytu nie jest kopią korutyny,
a zniszczenie uchwytu nie wywołuje automatycznie \texttt{destroy()}.
Dlatego wygodna klasa będąca właścicielem ramki powinna stosować RAII,
uniemożliwiać kopiowanie oraz poprawnie obsługiwać przenoszenie.

\begin{lstlisting}[caption={Operacje na uchwycie; fragment ilustracyjny.}]
#include <coroutine>

using Handle = std::coroutine_handle<>;

// Zakladamy, ze h wskazuje na istniejaca, zawieszona korutyne.
// if (h && !h.done()) h.resume();
// h.destroy(); // dokladnie raz, przez wlasciciela ramki
// Po destroy() wszystkie kopie h sa wiszacymi uchwytami.
\end{lstlisting}

Uchwyt może zostać przekazany do innego wątku lub executora, ale
wówczas trzeba zagwarantować, że ramka nie zostanie zniszczona,
zanim planista faktycznie wznowi zawieszoną korutynę. Jest to jedna
z głównych trudności implementacyjnych asynchronicznego API C++20.

\section{Własny generator: \texttt{co\_yield} od podstaw}
\label{sec:cpp15-generator}

Zbudujemy teraz typ \texttt{Generator}, który produkuje liczby całkowite
kolejno, \emph{na żądanie}. Poniższy przykład jest samowystarczalny.
Prywatny uchwyt do ramki stanowi jedynego właściciela zasobu;
\texttt{next()} wznawia korutynę tylko wtedy, gdy nie zakończyła
wykonywania. Własną obsługą wyjątku uzupełniamy przykład o
poprawną propagację błędów.

\begin{lstlisting}[caption={Generator liczb całkowitych: co\_yield i RAII (C++20); kompletny program.}]
#include <cassert>
#include <coroutine>
#include <exception>
#include <limits>
#include <stdexcept>
#include <utility>

class Generator {
public:
    struct promise_type {
        int current = 0;
        std::exception_ptr error;

        Generator get_return_object() {
            return Generator{Handle::from_promise(*this)};
        }
        std::suspend_always initial_suspend() noexcept { return {}; }
        std::suspend_always final_suspend() noexcept { return {}; }
        std::suspend_always yield_value(int x) noexcept {
            current = x;
            return {};
        }
        void return_void() noexcept {}
        void unhandled_exception() noexcept {
            error = std::current_exception();
        }
    };

    using Handle = std::coroutine_handle<promise_type>;

    explicit Generator(Handle h) noexcept : handle_(h) {}
    Generator(const Generator&) = delete;
    Generator& operator=(const Generator&) = delete;

    Generator(Generator&& other) noexcept
        : handle_(std::exchange(other.handle_, {})),
          ready_(std::exchange(other.ready_, false)) {}

    Generator& operator=(Generator&& other) noexcept {
        if (this != &other) {
            if (handle_) handle_.destroy();
            handle_ = std::exchange(other.handle_, {});
            ready_ = std::exchange(other.ready_, false);
        }
        return *this;
    }

    ~Generator() {
        if (handle_) handle_.destroy();
    }

    bool next() {
        ready_ = false;
        if (!handle_ || handle_.done()) return false;
        handle_.resume();
        if (handle_.promise().error)
            std::rethrow_exception(handle_.promise().error);
        ready_ = !handle_.done();
        return ready_;
    }

    int value() const {
        if (!ready_)
            throw std::logic_error("No current value");
        return handle_.promise().current;
    }

private:
    Handle handle_{};
    bool ready_ = false;
};

Generator squares(int n) {
    for (int i = 1; i <= n; ++i) {
        if (i > std::numeric_limits<int>::max() / i)
            throw std::overflow_error("square outside int");
        co_yield i * i;
    }
}

int main() {
    auto gen = squares(4); // Cialo jeszcze sie nie wykonalo.
    int sum = 0;
    while (gen.next()) sum += gen.value();
    assert(sum == 1 + 4 + 9 + 16);
    assert(!gen.next());
}
\end{lstlisting}

\subsection{Co dzieje się podczas \texttt{co\_yield}?}

Instrukcja \texttt{co\_yield i * i} powoduje wywołanie
\texttt{promise.yield\_value(i * i)} i zastosowanie mechanizmu
\texttt{co\_await} do otrzymanego awaitera. W naszym przypadku
\texttt{yield\_value} zapisuje wartość w obietnicy i zwraca
\texttt{suspend\_always}. Funkcja \texttt{next()} otrzymuje wtedy
sterowanie z powrotem i udostępnia zapisany element. Następne
\texttt{next()} kontynuuje pętlę od poprzedniego miejsca zawieszenia.

Generator posiada tylko jeden bieżący wynik. Nie przechowuje
całego ciągu jak \texttt{std::vector<int>}. To użyteczne przy dużych
lub potencjalnie nieskończonych sekwencjach. \emph{Nie oznacza to}
jednak automatycznie mniejszego kosztu obliczeń --- dochodzą koszty
ramki i kolejnych wznowień.

\nota{Przedstawiony \texttt{Generator} jest celowo minimalny: nie
udostępnia interfejsu iteratorów \texttt{begin()/end()}, nie jest
bezpieczny do jednoczesnego wznawiania z kilku wątków i nie powinien
być stosowany jako gotowa biblioteka produkcyjna.}

\section{Protokół \texttt{co\_await}: \texttt{await\_ready},
\texttt{await\_suspend}, \texttt{await\_resume}}
\label{sec:cpp15-awaiter}

Operator \texttt{co\_await} nie jest synonimem oczekiwania na
\texttt{future}. Kompilator uzyskuje obiekt \emph{awaiter}, którego
interfejs ma trzy fundamentalne operacje:
\begin{description}
\item[\texttt{await\_ready()}] Zwraca wartość logiczną: jeśli
      \texttt{true}, zawieszenie jest pomijane.
\item[\texttt{await\_suspend(handle)}] Jeżeli \texttt{await\_ready()}
      było fałszywe, metoda podejmuje decyzję o kontynuacji.
      Dopuszczalne postacie zwracają \texttt{void}, \texttt{bool}
      albo inny \texttt{std::coroutine\_handle}.
\item[\texttt{await\_resume()}] Wyznacza wynik całego wyrażenia
      \texttt{co\_await}; może też zgłosić wyjątek.
\end{description}

Gdy \texttt{await\_suspend} zwraca \texttt{bool}, wartość
\texttt{false} oznacza kontynuację bez pozostawania w zawieszeniu,
a \texttt{true} --- pozostawienie korutyny w stanie zawieszonym.
Przy zwrocie \texttt{void} korutyna pozostaje zawieszona do późniejszego
wznowienia. Zwrócenie innego uchwytu powoduje \emph{symmetric transfer}:
sterowanie przechodzi do wskazanej korutyny.

Nie każdy typ bezpośrednio dostarcza trzy wymienione metody: język
obsługuje też przeciążenie \texttt{operator co\_await} oraz
\texttt{await\_transform} zdefiniowane w \texttt{promise\_type}.
Sposób uzyskania awaitera zależy od tych reguł.

\subsection{Czysto lokalne zawieszenie bez nowego wątku}

Poniższy typ liczy wznowienia. Wykorzystamy go z prostą korutyną,
której właściciel ręcznie wywołuje \texttt{resume()}.

\begin{lstlisting}[caption={co\_await z własnym awaiterem: żadnego nowego wątku (C++20); kompletny program.}]
#include <cassert>
#include <coroutine>
#include <exception>
#include <utility>

struct Stepper {
    struct promise_type {
        Stepper get_return_object() {
            return Stepper{handle::from_promise(*this)};
        }
        std::suspend_always initial_suspend() noexcept { return {}; }
        std::suspend_always final_suspend() noexcept { return {}; }
        void return_void() noexcept {}
        void unhandled_exception() { std::terminate(); }
    };

    using handle = std::coroutine_handle<promise_type>;
    handle h;
    explicit Stepper(handle x) : h(x) {}
    Stepper(const Stepper&) = delete;
    Stepper(Stepper&& other) noexcept
        : h(std::exchange(other.h, {})) {}
    ~Stepper() { if (h) h.destroy(); }
    bool resume() {
        if (!h || h.done()) return false;
        h.resume();
        return !h.done();
    }
};

struct Pause {
    int& count;
    bool await_ready() const noexcept { return false; }
    void await_suspend(std::coroutine_handle<>) const noexcept {
        ++count; // Zawiesz, az wlasciciel wywola resume().
    }
    void await_resume() const noexcept {}
};

Stepper example(int& count) {
    ++count;                 // wykona sie przy pierwszym resume
    co_await Pause{count};  // pierwsze zawieszenie
    ++count;                 // wykona sie przy drugim resume
}

int main() {
    int count = 0;
    auto step = example(count);
    assert(count == 0);
    assert(step.resume() && count == 2);
    assert(!step.resume() && count == 3);
}
\end{lstlisting}

Kod niczego nie oblicza równolegle. Wszystkie instrukcje wykonał
wątek \texttt{main}. Korutyna jedynie oddała sterowanie, zachowując
miejsce wznowienia.

\section{Własne \texttt{Task<T>}: leniwe obliczenie a asynchroniczność}
\label{sec:cpp15-task}

Zaprojektujemy najmniejszy rozsądny typ \texttt{Task<T>} zwracający
wynik i obsługujący wyjątek. Celowo nie implementujemy jeszcze
harmonogramu: pozwala to zrozumieć, że \emph{sam typ korutyny nie jest
schedulerem}. \texttt{get()} w tym przykładzie działa wyłącznie
wtedy, gdy ciało może zakończyć się po pojedynczym wznowieniu.

\begin{lstlisting}[caption={Lazy Task<T> bez schedulera (C++20); kompletny program.}]
#include <cassert>
#include <coroutine>
#include <exception>
#include <optional>
#include <stdexcept>
#include <utility>

template<class T>
class Task {
public:
    struct promise_type {
        std::optional<T> result;
        std::exception_ptr error;
        Task get_return_object() {
            return Task{Handle::from_promise(*this)};
        }
        std::suspend_always initial_suspend() noexcept { return {}; }
        std::suspend_always final_suspend() noexcept { return {}; }
        void return_value(T x) { result.emplace(std::move(x)); }
        void unhandled_exception() noexcept {
            error = std::current_exception();
        }
    };
    using Handle = std::coroutine_handle<promise_type>;

    explicit Task(Handle h) noexcept : h_(h) {}
    Task(const Task&) = delete;
    Task& operator=(const Task&) = delete;
    Task(Task&& x) noexcept
        : h_(std::exchange(x.h_, {})),
          taken_(std::exchange(x.taken_, false)) {}
    ~Task() { if (h_) h_.destroy(); }

    T get() {
        if (!h_) throw std::logic_error("Empty task");
        if (taken_) throw std::logic_error("Result already consumed");
        if (!h_.done()) h_.resume();
        if (!h_.done())
            throw std::logic_error("Task needs an external scheduler");
        taken_ = true;
        if (h_.promise().error)
            std::rethrow_exception(h_.promise().error);
        return std::move(*h_.promise().result);
    }

private:
    Handle h_{};
    bool taken_ = false;
};

Task<int> calculate(int x) {
    co_return x * x;
}

int main() {
    auto task = calculate(7); // lazy, obliczenie jeszcze nie ruszylo
    assert(task.get() == 49);
}
\end{lstlisting}

Ten \texttt{Task<T>} jest pedagogiczny: nie ma przeciążenia
\texttt{operator co\_await}, nie wspiera wielokrotnego pobierania
wyniku (drugie \texttt{get()} zgłasza błąd) i nie potrafi samodzielnie
czekać na zdarzenia zewnętrzne.
W szczególności \texttt{get()} nie może po prostu wielokrotnie
wywoływać \texttt{resume()} na rzeczywistym zadaniu I/O: wznowienie
przed ukończeniem operacji mogłoby naruszyć protokół awaitera.
Produkcyjne \texttt{Task<T>} wymaga m.in. protokołu kontynuacji,
jasnej semantyki własności i poprawnej synchronizacji z executorem.

\section{Korutyna a \texttt{std::async}: fundamentalne porównanie}
\label{sec:cpp15-async-comparison}

\texttt{std::async} (C++11, nagłówek \texttt{<future>}) jest gotową
funkcją uruchamiającą zadanie i zwracającą \texttt{std::future<T>}.
Nie wymaga od użytkownika implementowania \texttt{promise\_type}.
Przy polityce \texttt{std::launch::async} funkcja pracuje w odrębnym
wątku; polityka \texttt{std::launch::deferred} odracza jej uruchomienie
do pierwszego nieczasowego oczekiwania na wynik.

Z kolei korutyna C++20 definiuje sposób zapisywania \emph{kontynuacji}.
Nawet jeśli wynik ma charakter asynchroniczny, \emph{faktyczne
uruchamianie pracy} zależy od awaitera, event loop, schedulera lub
innego mechanizmu zewnętrznego. Nie istnieje standardowa ogólna
reguła ,,jedna korutyna = jeden wątek''.


\paragraph{Czy można napisać \texttt{co\_await future}?}
Nie wprost: w C++20 standardowy \texttt{std::future<T>} nie udostępnia
protokołu awaitera wymaganego przez \texttt{co\_await}.
\texttt{future.get()} jest operacją potencjalnie blokującą
bieżący wątek. Istnieją biblioteki dostarczające awaitery dla
przyszłych wyników, ale sam adapter nie zamienia blokującej
operacji w nieblokującą. Naiwne uruchamianie osobnego wątku
tylko po to, by wykonał \texttt{future.wait()} i wznowił korutynę,
może zniweczyć zalety oszczędnego używania wątków.
Dla rzeczywiście nieblokującego oczekiwania na I/O potrzeba
integracji z właściwym mechanizmem asynchronicznym.

\begin{center}
\small
\begin{tabularx}{\textwidth}{@{}lXX@{}}
\toprule
\textbf{Cecha} & \textbf{\texttt{std::async}} & \textbf{Korutyny C++20} \\
\midrule
Co opisuje? & Zadanie ze stanem wyniku & Wykonanie z punktami zawieszenia \\
Tworzy nowy wątek? & \texttt{launch::async}: osobny wątek & Samo \texttt{co\_await}: nie \\
Kiedy zaczyna pracę? & Zależnie od polityki & Zależnie od \texttt{initial\_suspend} i schedulera \\
Odbieranie wyniku & \texttt{future::get()} (może blokować) & Zależy od typu \texttt{Task} i awaitera \\
Wyjątki & Przechowywane w stanie \texttt{future} & Obsługa przez \texttt{unhandled\_exception} i typ wyniku \\
Anulowanie & Brak automatycznego anulowania przez \texttt{future} & Możliwe kooperacyjnie; wymaga protokołu \\
Pamięć/zasoby & Stan asynchroniczny, przy \texttt{async} również wątek & Ramka korutyny, ewentualnie executor \\
Gotowe narzędzie w C++20? & Tak & Składnia tak; ogólny \texttt{Task<T>} nie \\
\bottomrule
\end{tabularx}
\end{center}

\subsection{Ten sam rachunek zapisany na dwa sposoby}

Porównanie dwóch krótkich funkcji pokazuje różnicę semantyczną.

\begin{lstlisting}[caption={std::async uruchamia zadanie w odrębnym wątku (C++11); kompletny program.}]
#include <cassert>
#include <future>

int square(int x) { return x * x; }

int main() {
    auto f = std::async(std::launch::async, square, 7);
    assert(f.get() == 49);
}
\end{lstlisting}

\begin{lstlisting}[caption={Korutyna bez awaitera zewnętrznego nie tworzy wątku; fragment.}]
// Przyjmujemy implementacje Task<T> z poprzedniej sekcji.
Task<int> square_coroutine(int x) {
    co_return x * x;
}

// Task<int> t = square_coroutine(7);
// int result = t.get(); // 49, w watku wywolujacym get()
\end{lstlisting}

W obu przypadkach można otrzymać tę samą wartość, ale nie jest to
to samo zachowanie wykonawcze. \texttt{std::async} z \texttt{launch::async}
zarządza wykonaniem w innym wątku, a pokazany \texttt{Task}
jedynie odkłada wywołanie do momentu ręcznego wznowienia.

\subsection{Pułapka destruktora \texttt{future}}

Ważna różnica dotyczy czasu życia. Gdy ostatni uchwyt do stanu
asynchronicznego utworzonego przez \texttt{std::async} z polityką
\texttt{launch::async} jest niszczony i praca wciąż trwa, destruktor
czeka na zakończenie powiązanego wątku.
Pierwsze z poniższych wywołań zakończy się, zanim rozpocznie się drugie:

\begin{lstlisting}[caption={Zignorowany future wymusza oczekiwanie; fragment.}]
// Funkcje slow_a i slow_b wykonuja dlugie obliczenia.
// std::async(std::launch::async, slow_a);
// std::async(std::launch::async, slow_b);
\end{lstlisting}

Jeżeli chcemy umożliwić nakładanie się zadań, trzeba przechować ich
uchwyty przed oczekiwaniem:
\begin{lstlisting}[caption={Dwa przyszłe wyniki powinny żyć jednocześnie; fragment.}]
// auto fa = std::async(std::launch::async, slow_a);
// auto fb = std::async(std::launch::async, slow_b);
// auto a = fa.get();
// auto b = fb.get();
\end{lstlisting}

Analogicznie w systemie korutyn należy znać politykę niszczenia
\texttt{Task}: usunięcie właściciela zawieszonej korutyny bez
koordynacji z executorem może zostawić wiszący uchwyt. Różne
biblioteki oferują odmienne kontrakty anulowania i oczekiwania.

\subsection{Asynchroniczne I/O a zadania CPU-bound}

Przy \emph{operacji I/O} większość czasu może upływać na oczekiwaniu
na dysk, sieć lub timer. Duża liczba korutyn może opisywać wiele
rozpoczętych operacji, podczas gdy mała pula wątków obsługuje
zdarzenia gotowości i wznowienia. Korzyść nie bierze się z tego,
że \texttt{co\_await} ,,jest szybsze'', lecz z możliwości
nieblokowania wątków podczas oczekiwania.

Przy \emph{operacjach CPU-bound} (obliczenia numeryczne, kompresja,
przetwarzanie obrazów) samo zawieszenie korutyny nie zwiększa
liczby rdzeni wykonujących kod. Tu potrzebne jest rzeczywiste
rozłożenie pracy na wiele wątków (np. \texttt{std::async}, pula
wątków lub OpenMP). Korutyna może być wygodnym \emph{interfejsem}
do takiego executora, ale nie zastępuje go.

\section{Rzeczywista asynchroniczność: korutyna i własny executor}
\label{sec:cpp15-executor}

Żeby zestawienie z \texttt{std::async} nie pozostało czysto
teoretyczne, utworzymy prosty executor z jednym wątkiem roboczym.
\texttt{co\_await executor.schedule()} odłoży kontynuację do
kolejki. Wznowienie wykona \emph{wątek executora}, nie wywołujący.
Wartością zwracaną będzie nasza klasa \texttt{FutureTask<T>},
wykorzystująca wewnętrznie \texttt{std::promise<T>} i
\texttt{std::future<T>}.

W przeciwieństwie do poprzedniego przykładu ramka ma strategię
\emph{samozniszczenia} przez \texttt{final\_suspend = suspend\_never}.
Nie przechowujemy już zewnętrznego uchwytu właścicielskiego:
po zakończeniu korutyna sama domyka swój cykl życia, a klient posiada
wyłącznie \texttt{future}. W tej uproszczonej konstrukcji executor
musi żyć aż do zakończenia zadań, a pojedyncza korutyna jest
wznawiana w kolejce co najwyżej raz.

\begin{lstlisting}[caption={Korutyna przechodząca do innego wątku poprzez executor (C++20); kompletny program.}]
#include <cassert>
#include <condition_variable>
#include <coroutine>
#include <deque>
#include <exception>
#include <future>
#include <mutex>
#include <stdexcept>
#include <thread>
#include <utility>

class Executor {
    std::mutex m_;
    std::condition_variable cv_;
    std::deque<std::coroutine_handle<>> queue_;
    bool closed_ = false;
    std::thread worker_;

    void run() {
        for (;;) {
            std::coroutine_handle<> h;
            {
                std::unique_lock lock(m_);
                cv_.wait(lock, [&] { return closed_ || !queue_.empty(); });
                if (queue_.empty() && closed_) return;
                h = queue_.front();
                queue_.pop_front();
            }
            h.resume(); // H moze sie tu samozniszczyc.
            // Nie korzystamy juz z h po resume().
        }
    }

public:
    Executor() : worker_([this] { run(); }) {}
    Executor(const Executor&) = delete;
    Executor& operator=(const Executor&) = delete;

    ~Executor() {
        {
            std::lock_guard lock(m_);
            closed_ = true;
        }
        cv_.notify_one();
        worker_.join();
    }

    void post(std::coroutine_handle<> h) {
        {
            std::lock_guard lock(m_);
            if (closed_) throw std::runtime_error("Executor closed");
            queue_.push_back(h);
        }
        cv_.notify_one();
    }

    struct Schedule {
        Executor& ex;
        bool await_ready() const noexcept { return false; }
        void await_suspend(std::coroutine_handle<> h) const {
            ex.post(h); // Po przekazaniu h nie uzywamy juz awaitera.
        }
        void await_resume() const noexcept {}
    };

    Schedule schedule() noexcept { return Schedule{*this}; }
};

template<class T>
class FutureTask {
    std::future<T> future_;
public:
    explicit FutureTask(std::future<T> f) : future_(std::move(f)) {}
    FutureTask(FutureTask&&) = default;
    FutureTask(const FutureTask&) = delete;

    T get() { return future_.get(); }

    struct promise_type {
        std::promise<T> result;
        FutureTask get_return_object() {
            return FutureTask{result.get_future()};
        }
        std::suspend_never initial_suspend() noexcept { return {}; }
        std::suspend_never final_suspend() noexcept { return {}; }
        void return_value(T x) { result.set_value(std::move(x)); }
        void unhandled_exception() noexcept {
            try { result.set_exception(std::current_exception()); }
            catch (...) { std::terminate(); }
        }
    };
};

FutureTask<int> square_on(Executor& ex, int x) {
    co_await ex.schedule(); // Tu oddajemy sterowanie.
    co_return x * x;        // Wykona sie w watku roboczym.
}

int main() {
    Executor ex;
    auto a = square_on(ex, 6);
    auto b = square_on(ex, 7);
    assert(a.get() + b.get() == 85);
    // Przed zniszczeniem ex oba zadania sie zakonczyly.
}
\end{lstlisting}

\paragraph{Co to porównanie rzeczywiście dowodzi?}
Kod z \texttt{square\_on} jest asynchroniczny względem wywołującego,
ponieważ po \texttt{co\_await} kontynuację przejmuje inny wątek.
Nie oznacza to jeszcze \emph{równoległego wykonywania} obu obliczeń:
executor ma tylko jednego pracownika, więc zadania są przez niego
wykonywane kolejno. Do równoległego przetwarzania dwóch obliczeń
należałoby dodać drugiego pracownika albo wykorzystać istniejącą
pulę wątków.

Pokazany executor jest \emph{szkicem edukacyjnym}. Nie obsługuje
ogólnego anulowania ani ponownego planowania tej samej korutyny,
nie rozwiązuje wszystkich problemów czasu życia i nie stanowi
gotowej biblioteki. W szczególności klient nie ma uchwytu pozwalającego
anulować i bezpiecznie zniszczyć zawieszone zadanie. W systemie
produkcyjnym własność korutyny, zamknięcie executora i czas życia
powiązanych zasobów muszą stanowić jeden, spójny kontrakt.

\section{Obsługa wyjątków, anulowanie i czas życia danych}
\label{sec:cpp15-errors}

\subsection{Wyjątki a \texttt{unhandled\_exception}}

Nieobsłużony wyjątek z ciała korutyny nie może zostać pozostawiony
bez decyzji. Kompilator wywołuje \texttt{promise.unhandled\_exception()}.
W generatorze zapisaliśmy \texttt{std::exception\_ptr}; przy kolejnym
\texttt{next()} błąd zostanie zgłoszony w kodzie klienta. W naszym
\texttt{FutureTask<T>} wyjątek trafia do stanu \texttt{std::future}
i ujawnia się podczas \texttt{get()}. Jest to analogiczne do
propagacji błędów w \texttt{std::async}, ale sam mechanizm wymagał
jawnej implementacji w \texttt{promise\_type}.

Należy też uwzględnić wyjątki zgłaszane przez operacje awaitera
oraz błędy alokacji ramki. Awaria podczas tworzenia ramki może
wystąpić jeszcze przed wykonaniem ciała korutyny. Typ obietnicy
może opcjonalnie udostępniać \texttt{get\_return\_object\_on\_allocation\_failure()}
razem z odpowiednią, niezgłaszającą wyjątku alokacją, ale nie jest to
mechanizm automatycznie dostępny każdemu typowi korutyny.

\subsection{Referencje przechowywane w ramce}

Argumenty referencyjne nie wydłużają czasu życia obiektów.
Jeżeli korutyna przechowuje referencję do zmiennej lokalnej, która
zniknie podczas zawieszenia, jej późniejsze wznowienie może prowadzić
do użycia nieistniejącego już obiektu.

\begin{lstlisting}[caption={Niebezpieczny czas życia referencji w korutynie; fragment koncepcyjny.}]
// Zakladamy lazy Task<int> z poprzedniej sekcji.
Task<int> read_later(const int& ref) {
    co_return ref;
}

// NIEPOPRAWNE:
// auto task = read_later(123); // 123 znika na koncu tej instrukcji
// int value = task.get();     // odczyt przez wiszaca referencje
\end{lstlisting}

Z reguły bezpieczniej przechowywać wewnątrz ramki wartości
przeznaczone do późniejszego użycia (np. obiekt przekazany przez
wartość), a nie referencje do tymczasowych obiektów.
Trzeba też zwracać uwagę na czas życia domknięć lambda:
korutyna będąca lambdą i wykorzystująca przechwycenia może
po zawieszeniu odwoływać się do już zniszczonego obiektu domknięcia.
Sam fakt przeniesienia wynikowej korutyny nie musi temu zapobiegać.

\subsection{Anulowanie jest kooperacyjne}

\texttt{std::stop\_token} (C++20) umożliwia \emph{żądanie}, a nie
wymuszenie zakończenia zadania. Korutyna może sprawdzać token
przed rozpoczęciem kolejnych faz albo korzystać z awaitera,
który potrafi odpowiednio reagować na anulowanie. Nie wynika
z tego jednak, że każde \texttt{co\_await} automatycznie obsługuje
\texttt{stop\_token}.

\begin{lstlisting}[caption={Kooperacyjne anulowanie korutyny: sprawdzenie żądania w punktach kontrolnych; fragment.}]
// Zakladamy Task<int> z poprzedniej sekcji.
#include <stop_token>

Task<int> calculate(std::stop_token token) {
    if (token.stop_requested()) co_return -1;
    // Pierwszy krok obliczen...
    if (token.stop_requested()) co_return -1;
    // Kolejny krok obliczen...
    co_return 42;
}
\end{lstlisting}

W powyższym przykładzie token jest przekazywany przez wartość.
Sam \texttt{stop\_token} nie steruje czasem życia ramki i nie
wznawia korutyny, która czeka na zdarzenie. Do pełnego anulowania
operacji I/O potrzeba jeszcze protokołu z usługą, która wykonuje
rzeczywistą operację, oraz reguł określających \emph{kto} zniszczy
korutynę i \emph{kiedy} wolno to zrobić.

\section{Pułapki synchronizacji i planowania}

\subsection{Równoczesne wznowienie tej samej korutyny}

Nie można traktować \texttt{coroutine\_handle::resume()} jak
operacji atomowej nadającej się do równoczesnego wywoływania przez
wiele wątków. Trzeba zapewnić, że dana zawieszona korutyna ma
jednego logicznego właściciela uprawnienia do wznowienia.
Ponowne wznowienie korutyny już aktywnej lub zakończonej może
prowadzić do niezdefiniowanego zachowania.

\subsection{Utracone powiadomienie i wyścig publikacji uchwytu}

Awaiter może przekazać uchwyt do innego wątku wewnątrz
\texttt{await\_suspend}. Wówczas druga strona może wznowić korutynę
\emph{zanim} pierwotne wywołanie \texttt{await\_suspend} powróci.
Dlatego po opublikowaniu uchwytu kod awaitera nie powinien korzystać
z obiektów, których czas życia mógł się już zakończyć.

W bardziej złożonych awaiterach trzeba również uniknąć zgubienia
zdarzenia kończącego operację między sprawdzeniem gotowości a
zarejestrowaniem oczekiwania. Poprawne rozwiązania stosują np.
mutex, operacje atomowe albo protokół jednorazowego przejęcia
kontynuacji. Samo \texttt{await\_ready()} nie zastępuje synchronizacji.

\subsection{Brak własności \texttt{coroutine\_handle}}

Gdy planista ma w kolejce jedynie \texttt{coroutine\_handle<>},
nie posiada informacji, kto może zniszczyć ramkę. Usunięcie
obiektu \texttt{Task} podczas oczekiwania może prowadzić do
późniejszego wznowienia już zwolnionej pamięci.
Najbezpieczniejszy projekt definiuje własność od początku:
\emph{Task posiada ramkę}, \emph{executor posiada ramkę} albo
\emph{ramka posiada samą siebie}, z dokładnym protokołem przekazania
własności i zakończenia.

\section{Koszty, optymalizacje i pomiary}
\label{sec:cpp15-performance}

Wydajność należy rozpatrywać w całym systemie, a nie tylko przez
porównanie czasu jednego \texttt{co\_await} i jednego \texttt{std::async}.
\begin{itemize}
\item \textbf{Utworzenie wątku.} \texttt{std::async(launch::async, ...)}
      uruchamia pracę w odrębnym wątku; bardzo duża liczba krótkich
      zadań może generować istotny narzut zarządzania wątkami.
\item \textbf{Ramka.} Korutyna może wymagać alokacji pamięci,
      zachowywania stanu i wznawiania. Kompilator może czasem
      wyeliminować alokację ramki, ale nie jest to gwarancja.
\item \textbf{Executor.} Kolejka, mutexy, przełączanie między wątkami
      oraz współdzielenie pamięci również kosztują.
\item \textbf{Lokalność pamięci.} Duża liczba korutyn z rozproszonymi
      ramkami może prowadzić do dodatkowych nietrafień cache.
\item \textbf{Skalowalność.} Dla I/O ważniejsza może być liczba
      oczekujących operacji niż przepustowość pojedynczego zadania.
      Dla CPU-bound kluczowa jest liczba rdzeni, wielkość porcji pracy
      i ograniczenie dostępu do współdzielonych zasobów.
\end{itemize}

Rzetelny benchmark powinien mierzyć co najmniej: całkowitą
przepustowość, opóźnienia (również wartości odstające), liczbę wątków,
liczbę alokacji oraz maksymalne zużycie pamięci. Porównywanie
\texttt{std::async} z \emph{nieuruchamiającą nowych wątków} korutyną
wykonującą natychmiast prostą arytmetykę nie jest miarodajnym
testem wydajności asynchronicznej.

\subsection{Praktyczna decyzja: kiedy używać którego narzędzia?}

\begin{center}
\small
\begin{tabularx}{\textwidth}{@{}lX@{}}
\toprule
\textbf{Problem} & \textbf{Rozsądny pierwszy wybór} \\
\midrule
Jedno niezależne zadanie CPU z wynikiem & \texttt{std::async(std::launch::async, ...)} \\
Wiele krótkich zadań CPU & Pula wątków, executor lub OpenMP \\
Wiele połączeń I/O z długim oczekiwaniem & Korutyny + asynchroniczny runtime I/O \\
Leniwe generowanie sekwencji & Generator korutynowy \\
Wspólny licznik lub krótka sekcja krytyczna & \texttt{atomic} albo mutex, nie korutyna \\
Anulowalne zadanie działające długo & Przemyślany protokół \texttt{stop\_token} \\
Interfejs biblioteki o złożonych zależnościach & \texttt{Task<T>} nad zdefiniowanym executorem \\
\bottomrule
\end{tabularx}
\end{center}

\section{Związek z OpenMP, \texttt{thread} i \texttt{future}}

Rozszerzony rozdział o współbieżności omawiał \texttt{std::thread},
\texttt{std::jthread}, atomiki, blokady, \texttt{std::async} oraz OpenMP.
Korutyny nie unieważniają tych narzędzi; stanowią kolejną warstwę.
\begin{itemize}
\item \texttt{std::thread} opisuje konkretny wątek wykonania,
      a \texttt{std::jthread} dodatkowo pomaga zarządzać jego
      czasem życia i kooperacyjnym zatrzymywaniem.
\item \texttt{std::async} tworzy stan przyszłego wyniku i
      uruchamia wywołanie zgodnie z przyjętą polityką.
\item OpenMP rozdziela pętle i zadania między wątki obsługiwane
      przez środowisko wykonawcze OpenMP; szczególnie przydaje się
      w obliczeniach CPU-bound i kodzie numerycznym.
\item Korutyny pozwalają wygodnie opisać wieloetapową logikę
      oczekiwania i kontynuacji. Ich wydajność i stopień równoległości
      zależą od dostarczonego środowiska wykonawczego.
\end{itemize}

Możliwe jest użycie \emph{korutyn nad pulą wątków}: \texttt{co\_await}
przekazuje kontynuację do kolejki, a pracownik puli później ją wznawia.
Nie należy jednak od razu łączyć wszystkich mechanizmów.
Najpierw określamy, czy program czeka głównie na I/O, czy wykonuje
obliczenia, a następnie wybieramy model wykonania i własności zasobów.

\section{Najczęstsze błędy projektowe}

\begin{enumerate}
\item Zakładanie, że \texttt{co\_await} uruchamia nowy wątek.
\item Próba zwrócenia bezpośrednio \texttt{std::future<T>}
      z własnej korutyny bez odpowiedniego \texttt{promise\_type}.
\item Dwukrotne wywołanie \texttt{destroy()} dla tej samej ramki.
\item Wywołanie \texttt{resume()} po osiągnięciu
      \texttt{final\_suspend}.
\item Pozostawienie referencji lub przechwycenia lambdy do obiektu,
      który znika podczas zawieszenia.
\item Zniszczenie \texttt{Task} mimo tego, że executor posiada
      w kolejce uchwyt do jego ramki.
\item Mylenie automatycznej propagacji wyjątków w \texttt{std::async}
      z obowiązkami implementatora \texttt{promise\_type}.
\item Traktowanie \texttt{std::stop\_token} jako mechanizmu
      przymusowego przerwania dowolnego \texttt{co\_await}.
\item Równoczesne wznawianie tej samej korutyny z kilku wątków.
\item Uznanie, że zawieszenie i wznowienie zawsze kosztuje mniej
      niż zwykłe wywołanie funkcji.
\item Uruchamianie \texttt{std::async} dla tysięcy mikro-zadań
      bez rozważenia puli wątków.
\item Zwracanie z \texttt{await\_suspend} po opublikowaniu uchwytu
      i dalsze używanie potencjalnie zniszczonego awaitera.
\end{enumerate}

\section*{Podsumowanie rozdziału}
\addcontentsline{toc}{section}{Podsumowanie rozdziału}

Korutyny C++20 są mechanizmem \emph{zawieszenia i kontynuacji},
a nie samodzielnym środowiskiem asynchronicznym. Generator używa
\texttt{co\_yield}, aby leniwie udostępniać wartości. Typ zadania
używa \texttt{co\_return}, obiektu \texttt{promise\_type} i awaiterów,
aby opisać wykonanie z wynikiem i wyjątkami. \texttt{co\_await}
przekazuje kontrolę zgodnie z protokołem awaitera, a dopiero executor
lub inne środowisko może wznowić kod w osobnym wątku.

\texttt{std::async} jest prostym narzędziem do uruchomienia funkcji
i odebrania przyszłego wyniku. Korutyny są lepszym budulcem dla
złożonych przepływów wieloetapowych, szczególnie I/O, ale wymagają
większej odpowiedzialności za pamięć, synchronizację, anulowanie
i czas życia ramki. W obu przypadkach należy rozdzielać trzy pytania:
\emph{kto wykonuje kod, kto posiada stan i kto decyduje o zakończeniu}.

\section*{Literatura i dokumentacja}
\addcontentsline{toc}{section}{Literatura i dokumentacja}
\begin{itemize}
\item \url{https://en.cppreference.com/w/cpp/language/coroutines} --- model korutyn C++20.
\item \url{https://en.cppreference.com/w/cpp/coroutine/coroutine_handle} --- uchwyty i zarządzanie ramką.
\item \url{https://en.cppreference.com/w/cpp/coroutine/suspend_always} --- podstawowe awaitery.
\item \url{https://en.cppreference.com/w/cpp/thread/async} --- \texttt{std::async} i polityki uruchamiania.
\item \url{https://en.cppreference.com/w/cpp/thread/future} --- \texttt{std::future}.
\item \url{https://en.cppreference.com/w/cpp/coroutine/generator} --- \texttt{std::generator} (C++23).
\item Lewis Baker, \emph{C++ Coroutines: Understanding the Compiler Transformations}
      --- seria publikacji o modelu wykonania korutyn.
\item David Mazières, \emph{My Tutorial and Take on C++20 Coroutines}, 2021.
\end{itemize}

% Listing oznaczony jako „kompletny program” mozna skompilowac oddzielnie.

\chapter{Korutyny, zadania asynchroniczne i \texttt{std::async}}
\label{chap:cpp15-coroutines}

\section{Po co C++ potrzebuje korutyn?}

Zwykła funkcja ma zasadniczo jeden punkt wejścia. Może wywoływać inne
funkcje, ale po powrocie oddaje sterowanie wywołującemu i nie posiada
mechanizmu późniejszego wznowienia dokładnie w miejscu, w którym
zakończyła wykonywanie. \emph{Korutyna} (ang. \emph{coroutine}) jest
funkcją, której wykonanie można \emph{zawiesić}, a następnie
\emph{wznowić}, zachowując wybrane zmienne lokalne i miejsce kontynuacji.

Korutyny C++20 wspierają między innymi:
\begin{itemize}
\item \emph{generatory} --- wytwarzanie kolejnych wartości na żądanie,
      bez budowania całej kolekcji;
\item \emph{zadania asynchroniczne} --- opis operacji, których kolejne
      etapy uruchamiane są po zakończeniu wcześniejszych;
\item \emph{kooperacyjne planowanie} --- oddawanie sterowania
      harmonogramowi zadań;
\item \emph{maszyny stanów} --- zachowywanie stanu pomiędzy wznowieniami
      bez ręcznego kodowania każdej możliwej kontynuacji.
\end{itemize}

Najważniejsza zasada brzmi:
\begin{center}
\emph{Korutyna nie jest wątkiem, a \texttt{co\_await} nie oznacza
automatycznie ,,uruchom w nowym wątku''.}
\end{center}

Korutyna określa \emph{jak zatrzymać i kontynuować obliczenie};
\texttt{std::async}, \texttt{std::thread} i pula wątków dotyczą również
\emph{gdzie i kiedy} wykonywać pracę. W szczególności kod korutyny
może działać od początku do końca w jednym wątku, nawet jeśli
wielokrotnie się zawiesza. Odwrotnie, korutyna może być wznawiana
przez różne wątki, jeżeli tak zdecyduje dostarczony mechanizm planowania.

\subsection{Trzy słowa kluczowe}

\begin{center}\small
\begin{tabularx}{\textwidth}{@{}p{0.32\textwidth}X@{}}
\toprule
Konstrukcja & Działanie w korutynie \\
\midrule
\code{co_await expr} & Przekazuje sterowanie obiektowi oczekującemu; może zawiesić korutynę bez blokowania wątku. \\
\code{co_yield value} & Udostępnia kolejną wartość i zwykle zawiesza generator do wznowienia. \\
\code{co_return value} & Kończy korutynę i przekazuje wynik zgodnie z protokołem jej typu zwracanego. \\
\code{await_ready()} & Określa, czy wynik jest już gotowy i można pominąć zawieszenie. \\
\code{await_suspend(h)} & Obsługuje zawieszenie i uchwyt \code{h}; może zaplanować późniejsze wznowienie. \\
\code{await_resume()} & Zwraca rezultat oczekiwania albo zgłasza zapisany wyjątek po wznowieniu. \\
\bottomrule
\end{tabularx}
\end{center}

C++20 wprowadził następujące konstrukcje składniowe:
\begin{description}
\item[\texttt{co\_await expr}] Zastosuj protokół oczekiwania do wyrażenia;
      zależnie od niego możesz zawiesić korutynę i później ją wznowić.
\item[\texttt{co\_yield expr}] Udostępnij kolejną wartość generatora;
      składniowo jest to szczególna postać oczekiwania, wykorzystująca
      metodę \texttt{yield\_value} obiektu \texttt{promise\_type}.
\item[\texttt{co\_return expr}] Zakończ korutynę i przekaż wynik
      przez odpowiednie \texttt{return\_value} albo \texttt{return\_void}.
\end{description}

Funkcja zawierająca jeden z tych operatorów (w odpowiednim kontekście)
rozpoznawana jest jako korutyna. Typ jej zwracanej wartości musi
udostępniać właściwy \emph{typ obietnicy} \texttt{promise\_type}
(lub odpowiednią specjalizację \texttt{std::coroutine\_traits}).
To wymóg mechanizmu korutyn, a nie znanego już \texttt{std::promise}.

\begin{lstlisting}[caption={Niepoprawna intuicja: std::future nie jest automatycznie typem korutyny (C++20); fragment celowo błędny.}]
#include <future>

// Standardowe std::future<int> nie zapewnia promise_type:
// std::future<int> calculate() { co_return 42; } // BLAD!
\end{lstlisting}

W C++20 biblioteka standardowa dostarcza \emph{mechanizm} korutyn
(nagłówek \texttt{<coroutine>}), ale nie dostarcza gotowego, ogólnego
\texttt{Task<T>} ani standardowego executora. Takie typy definiuje
biblioteka lub programista. \texttt{std::generator} jest rozwiązaniem
późniejszym, wprowadzonym dopiero w C++23, a jego wsparcie zależy
od użytej biblioteki standardowej.

\section{Różnica między zawieszeniem a zablokowaniem wątku}
\label{sec:cpp15-suspend-block}

Rozważmy dwa sposoby wstrzymania działania programu:
\begin{itemize}
\item \texttt{std::this\_thread::sleep\_for(100ms)} \emph{blokuje
      bieżący wątek} na pewien czas; sam wątek pozostaje zajętym
      zasobem systemowym, choć może nie zużywać wówczas czasu CPU;
\item dobrze zaprojektowane \texttt{co\_await} \emph{zawiesza
      korutynę}, oddając wątek do wykonywania innych zadań;
      wznowienie może nastąpić po nadejściu wyniku operacji I/O,
      komunikatu lub upływie czasu.
\end{itemize}

Nie każde \texttt{co\_await} daje zawieszenie: awaiter może zgłosić,
że wynik jest już gotowy. Również samo użycie \texttt{co\_await}
nie gwarantuje nieblokującego I/O --- zależy to od jego implementacji.

\begin{center}
\small
\begin{tabularx}{\textwidth}{@{}lXX@{}}
\toprule
\textbf{Właściwość} & \textbf{Blokada wątku} & \textbf{Zawieszenie korutyny} \\
\midrule
Zachowanie stanu & Na stosie wywołań wątku & Zwykle w ramce korutyny \\
Wątek potrzebny podczas oczekiwania & Tak, istnieje i czeka & Nie jest konieczny \emph{dla samego oczekiwania} \\
Wznowienie & Planista systemu/odblokowanie & Mechanizm awaitera i właściciel kontynuacji \\
Koszt & M.in. wątek i jego stos & M.in. ramka i obsługa zawieszenia \\
\bottomrule
\end{tabularx}
\end{center}

\section{Ramka korutyny i jej cykl życia}
\label{sec:cpp15-frame}

Kompilator przekształca korutynę w obiekt reprezentujący jej stan.
Ten stan nazywamy \emph{ramką korutyny} (ang. \emph{coroutine frame}).
Może ona obejmować obiekt \texttt{promise\_type}, parametry i stan
potrzebny do kontynuowania obliczenia po zawieszeniu. To, które
zmienne lokalne trzeba zachować, wynika z przebiegu programu.

W ogólnym modelu wywołanie korutyny przebiega w następujących etapach:
\begin{enumerate}
\item Alokowana jest ramka, chyba że implementacja legalnie
      wyeliminuje tę alokację (np. przez zagnieżdżenie ramki w innym obiekcie).
\item Kopiowane lub przenoszone są parametry odpowiednio do ich rodzajów;
      \emph{parametry referencyjne pozostają referencjami}.
\item Konstruowany jest obiekt obietnicy \texttt{promise\_type}.
\item Wywoływana jest metoda \texttt{get\_return\_object()}.
\item Wykonywany jest \texttt{co\_await promise.initial\_suspend()}.
\item Uruchamiane jest ciało korutyny (od razu lub po wznowieniu).
\item \texttt{co\_return} uruchamia \texttt{return\_value} lub
      \texttt{return\_void}; nieobsłużony wyjątek trafia do
      \texttt{unhandled\_exception()}.
\item Niszczone są obiekty lokalne, a korutyna wykonuje końcowy
      \texttt{co\_await} na wyniku \texttt{promise.final\_suspend()}.
\item Ramka zostaje ostatecznie zniszczona we właściwym momencie.
\end{enumerate}

Kolejność i czas \emph{zwolnienia ramki} zależą od przyjętej polityki
własności. Dla \texttt{final\_suspend = suspend\_always} ramkę zwykle
niszczy właściciel uchwytu, nawet po zakończeniu ciała korutyny.
Przy \texttt{final\_suspend = suspend\_never} mechanizm korutyny
może zakończyć się bez pozostawiania ramki zawieszonej do późniejszego
zniszczenia przez właściciela. Tych strategii nie wolno mieszać.

\subsection{Początek i koniec: \texttt{initial\_suspend} i \texttt{final\_suspend}}

Dwa najczęściej używane typy z nagłówka \texttt{<coroutine>} to:
\begin{description}
\item[\texttt{std::suspend\_never}] Oczekiwanie jest natychmiast
      gotowe i nie powoduje zawieszenia.
\item[\texttt{std::suspend\_always}] Oczekiwanie zawsze prowadzi
      do zawieszenia, dopóki ktoś nie wznowi korutyny.
\end{description}

Jeżeli \texttt{initial\_suspend()} zwraca \texttt{suspend\_always},
korutyna jest \emph{lazy}: nie rozpoczyna wykonywania ciała przy samym
wywołaniu funkcji. Natomiast \texttt{suspend\_never} oznacza
\emph{eager start}: ciało może rozpocząć się natychmiast w wątku
wywołującym i pracować aż do pierwszego rzeczywistego zawieszenia.

\textbf{Ważne:} nie wolno wywoływać \texttt{resume()} na korutynie,
która jest zawieszona w swoim \texttt{final\_suspend}. Jest to
niezdefiniowane zachowanie. Uchwyt do takiej ramki może nadal służyć
do \texttt{destroy()}, jeżeli ramka nie została już zniszczona.

\section{Uchwyt \texttt{std::coroutine\_handle} i własność ramki}

\texttt{std::coroutine\_handle<Promise>} jest lekkim uchwytem
umożliwiającym m.in. wznowienie korutyny i dostęp do jej obietnicy.
Sam \texttt{coroutine\_handle} \emph{nie zarządza automatycznie
czasem życia ramki}. Skopiowanie uchwytu nie jest kopią korutyny,
a zniszczenie uchwytu nie wywołuje automatycznie \texttt{destroy()}.
Dlatego wygodna klasa będąca właścicielem ramki powinna stosować RAII,
uniemożliwiać kopiowanie oraz poprawnie obsługiwać przenoszenie.

\begin{lstlisting}[caption={Operacje na uchwycie; fragment ilustracyjny.}]
#include <coroutine>

using Handle = std::coroutine_handle<>;

// Zakladamy, ze h wskazuje na istniejaca, zawieszona korutyne.
// if (h && !h.done()) h.resume();
// h.destroy(); // dokladnie raz, przez wlasciciela ramki
// Po destroy() wszystkie kopie h sa wiszacymi uchwytami.
\end{lstlisting}

Uchwyt może zostać przekazany do innego wątku lub executora, ale
wówczas trzeba zagwarantować, że ramka nie zostanie zniszczona,
zanim planista faktycznie wznowi zawieszoną korutynę. Jest to jedna
z głównych trudności implementacyjnych asynchronicznego API C++20.

\section{Własny generator: \texttt{co\_yield} od podstaw}
\label{sec:cpp15-generator}

Zbudujemy teraz typ \texttt{Generator}, który produkuje liczby całkowite
kolejno, \emph{na żądanie}. Poniższy przykład jest samowystarczalny.
Prywatny uchwyt do ramki stanowi jedynego właściciela zasobu;
\texttt{next()} wznawia korutynę tylko wtedy, gdy nie zakończyła
wykonywania. Własną obsługą wyjątku uzupełniamy przykład o
poprawną propagację błędów.

\begin{lstlisting}[caption={Generator liczb całkowitych: co\_yield i RAII (C++20); kompletny program.}]
#include <cassert>
#include <coroutine>
#include <exception>
#include <limits>
#include <stdexcept>
#include <utility>

class Generator {
public:
    struct promise_type {
        int current = 0;
        std::exception_ptr error;

        Generator get_return_object() {
            return Generator{Handle::from_promise(*this)};
        }
        std::suspend_always initial_suspend() noexcept { return {}; }
        std::suspend_always final_suspend() noexcept { return {}; }
        std::suspend_always yield_value(int x) noexcept {
            current = x;
            return {};
        }
        void return_void() noexcept {}
        void unhandled_exception() noexcept {
            error = std::current_exception();
        }
    };

    using Handle = std::coroutine_handle<promise_type>;

    explicit Generator(Handle h) noexcept : handle_(h) {}
    Generator(const Generator&) = delete;
    Generator& operator=(const Generator&) = delete;

    Generator(Generator&& other) noexcept
        : handle_(std::exchange(other.handle_, {})),
          ready_(std::exchange(other.ready_, false)) {}

    Generator& operator=(Generator&& other) noexcept {
        if (this != &other) {
            if (handle_) handle_.destroy();
            handle_ = std::exchange(other.handle_, {});
            ready_ = std::exchange(other.ready_, false);
        }
        return *this;
    }

    ~Generator() {
        if (handle_) handle_.destroy();
    }

    bool next() {
        ready_ = false;
        if (!handle_ || handle_.done()) return false;
        handle_.resume();
        if (handle_.promise().error)
            std::rethrow_exception(handle_.promise().error);
        ready_ = !handle_.done();
        return ready_;
    }

    int value() const {
        if (!ready_)
            throw std::logic_error("No current value");
        return handle_.promise().current;
    }

private:
    Handle handle_{};
    bool ready_ = false;
};

Generator squares(int n) {
    for (int i = 1; i <= n; ++i) {
        if (i > std::numeric_limits<int>::max() / i)
            throw std::overflow_error("square outside int");
        co_yield i * i;
    }
}

int main() {
    auto gen = squares(4); // Cialo jeszcze sie nie wykonalo.
    int sum = 0;
    while (gen.next()) sum += gen.value();
    assert(sum == 1 + 4 + 9 + 16);
    assert(!gen.next());
}
\end{lstlisting}

\subsection{Co dzieje się podczas \texttt{co\_yield}?}

Instrukcja \texttt{co\_yield i * i} powoduje wywołanie
\texttt{promise.yield\_value(i * i)} i zastosowanie mechanizmu
\texttt{co\_await} do otrzymanego awaitera. W naszym przypadku
\texttt{yield\_value} zapisuje wartość w obietnicy i zwraca
\texttt{suspend\_always}. Funkcja \texttt{next()} otrzymuje wtedy
sterowanie z powrotem i udostępnia zapisany element. Następne
\texttt{next()} kontynuuje pętlę od poprzedniego miejsca zawieszenia.

Generator posiada tylko jeden bieżący wynik. Nie przechowuje
całego ciągu jak \texttt{std::vector<int>}. To użyteczne przy dużych
lub potencjalnie nieskończonych sekwencjach. \emph{Nie oznacza to}
jednak automatycznie mniejszego kosztu obliczeń --- dochodzą koszty
ramki i kolejnych wznowień.

\nota{Przedstawiony \texttt{Generator} jest celowo minimalny: nie
udostępnia interfejsu iteratorów \texttt{begin()/end()}, nie jest
bezpieczny do jednoczesnego wznawiania z kilku wątków i nie powinien
być stosowany jako gotowa biblioteka produkcyjna.}

\section{Protokół \texttt{co\_await}: \texttt{await\_ready},
\texttt{await\_suspend}, \texttt{await\_resume}}
\label{sec:cpp15-awaiter}

Operator \texttt{co\_await} nie jest synonimem oczekiwania na
\texttt{future}. Kompilator uzyskuje obiekt \emph{awaiter}, którego
interfejs ma trzy fundamentalne operacje:
\begin{description}
\item[\texttt{await\_ready()}] Zwraca wartość logiczną: jeśli
      \texttt{true}, zawieszenie jest pomijane.
\item[\texttt{await\_suspend(handle)}] Jeżeli \texttt{await\_ready()}
      było fałszywe, metoda podejmuje decyzję o kontynuacji.
      Dopuszczalne postacie zwracają \texttt{void}, \texttt{bool}
      albo inny \texttt{std::coroutine\_handle}.
\item[\texttt{await\_resume()}] Wyznacza wynik całego wyrażenia
      \texttt{co\_await}; może też zgłosić wyjątek.
\end{description}

Gdy \texttt{await\_suspend} zwraca \texttt{bool}, wartość
\texttt{false} oznacza kontynuację bez pozostawania w zawieszeniu,
a \texttt{true} --- pozostawienie korutyny w stanie zawieszonym.
Przy zwrocie \texttt{void} korutyna pozostaje zawieszona do późniejszego
wznowienia. Zwrócenie innego uchwytu powoduje \emph{symmetric transfer}:
sterowanie przechodzi do wskazanej korutyny.

Nie każdy typ bezpośrednio dostarcza trzy wymienione metody: język
obsługuje też przeciążenie \texttt{operator co\_await} oraz
\texttt{await\_transform} zdefiniowane w \texttt{promise\_type}.
Sposób uzyskania awaitera zależy od tych reguł.

\subsection{Czysto lokalne zawieszenie bez nowego wątku}

Poniższy typ liczy wznowienia. Wykorzystamy go z prostą korutyną,
której właściciel ręcznie wywołuje \texttt{resume()}.

\begin{lstlisting}[caption={co\_await z własnym awaiterem: żadnego nowego wątku (C++20); kompletny program.}]
#include <cassert>
#include <coroutine>
#include <exception>
#include <utility>

struct Stepper {
    struct promise_type {
        Stepper get_return_object() {
            return Stepper{handle::from_promise(*this)};
        }
        std::suspend_always initial_suspend() noexcept { return {}; }
        std::suspend_always final_suspend() noexcept { return {}; }
        void return_void() noexcept {}
        void unhandled_exception() { std::terminate(); }
    };

    using handle = std::coroutine_handle<promise_type>;
    handle h;
    explicit Stepper(handle x) : h(x) {}
    Stepper(const Stepper&) = delete;
    Stepper(Stepper&& other) noexcept
        : h(std::exchange(other.h, {})) {}
    ~Stepper() { if (h) h.destroy(); }
    bool resume() {
        if (!h || h.done()) return false;
        h.resume();
        return !h.done();
    }
};

struct Pause {
    int& count;
    bool await_ready() const noexcept { return false; }
    void await_suspend(std::coroutine_handle<>) const noexcept {
        ++count; // Zawiesz, az wlasciciel wywola resume().
    }
    void await_resume() const noexcept {}
};

Stepper example(int& count) {
    ++count;                 // wykona sie przy pierwszym resume
    co_await Pause{count};  // pierwsze zawieszenie
    ++count;                 // wykona sie przy drugim resume
}

int main() {
    int count = 0;
    auto step = example(count);
    assert(count == 0);
    assert(step.resume() && count == 2);
    assert(!step.resume() && count == 3);
}
\end{lstlisting}

Kod niczego nie oblicza równolegle. Wszystkie instrukcje wykonał
wątek \texttt{main}. Korutyna jedynie oddała sterowanie, zachowując
miejsce wznowienia.

\section{Własne \texttt{Task<T>}: leniwe obliczenie a asynchroniczność}
\label{sec:cpp15-task}

Zaprojektujemy najmniejszy rozsądny typ \texttt{Task<T>} zwracający
wynik i obsługujący wyjątek. Celowo nie implementujemy jeszcze
harmonogramu: pozwala to zrozumieć, że \emph{sam typ korutyny nie jest
schedulerem}. \texttt{get()} w tym przykładzie działa wyłącznie
wtedy, gdy ciało może zakończyć się po pojedynczym wznowieniu.

\begin{lstlisting}[caption={Lazy Task<T> bez schedulera (C++20); kompletny program.}]
#include <cassert>
#include <coroutine>
#include <exception>
#include <optional>
#include <stdexcept>
#include <utility>

template<class T>
class Task {
public:
    struct promise_type {
        std::optional<T> result;
        std::exception_ptr error;
        Task get_return_object() {
            return Task{Handle::from_promise(*this)};
        }
        std::suspend_always initial_suspend() noexcept { return {}; }
        std::suspend_always final_suspend() noexcept { return {}; }
        void return_value(T x) { result.emplace(std::move(x)); }
        void unhandled_exception() noexcept {
            error = std::current_exception();
        }
    };
    using Handle = std::coroutine_handle<promise_type>;

    explicit Task(Handle h) noexcept : h_(h) {}
    Task(const Task&) = delete;
    Task& operator=(const Task&) = delete;
    Task(Task&& x) noexcept
        : h_(std::exchange(x.h_, {})),
          taken_(std::exchange(x.taken_, false)) {}
    ~Task() { if (h_) h_.destroy(); }

    T get() {
        if (!h_) throw std::logic_error("Empty task");
        if (taken_) throw std::logic_error("Result already consumed");
        if (!h_.done()) h_.resume();
        if (!h_.done())
            throw std::logic_error("Task needs an external scheduler");
        taken_ = true;
        if (h_.promise().error)
            std::rethrow_exception(h_.promise().error);
        return std::move(*h_.promise().result);
    }

private:
    Handle h_{};
    bool taken_ = false;
};

Task<int> calculate(int x) {
    co_return x * x;
}

int main() {
    auto task = calculate(7); // lazy, obliczenie jeszcze nie ruszylo
    assert(task.get() == 49);
}
\end{lstlisting}

Ten \texttt{Task<T>} jest pedagogiczny: nie ma przeciążenia
\texttt{operator co\_await}, nie wspiera wielokrotnego pobierania
wyniku (drugie \texttt{get()} zgłasza błąd) i nie potrafi samodzielnie
czekać na zdarzenia zewnętrzne.
W szczególności \texttt{get()} nie może po prostu wielokrotnie
wywoływać \texttt{resume()} na rzeczywistym zadaniu I/O: wznowienie
przed ukończeniem operacji mogłoby naruszyć protokół awaitera.
Produkcyjne \texttt{Task<T>} wymaga m.in. protokołu kontynuacji,
jasnej semantyki własności i poprawnej synchronizacji z executorem.

\section{Korutyna a \texttt{std::async}: fundamentalne porównanie}
\label{sec:cpp15-async-comparison}

\texttt{std::async} (C++11, nagłówek \texttt{<future>}) jest gotową
funkcją uruchamiającą zadanie i zwracającą \texttt{std::future<T>}.
Nie wymaga od użytkownika implementowania \texttt{promise\_type}.
Przy polityce \texttt{std::launch::async} funkcja pracuje w odrębnym
wątku; polityka \texttt{std::launch::deferred} odracza jej uruchomienie
do pierwszego nieczasowego oczekiwania na wynik.

Z kolei korutyna C++20 definiuje sposób zapisywania \emph{kontynuacji}.
Nawet jeśli wynik ma charakter asynchroniczny, \emph{faktyczne
uruchamianie pracy} zależy od awaitera, event loop, schedulera lub
innego mechanizmu zewnętrznego. Nie istnieje standardowa ogólna
reguła ,,jedna korutyna = jeden wątek''.


\paragraph{Czy można napisać \texttt{co\_await future}?}
Nie wprost: w C++20 standardowy \texttt{std::future<T>} nie udostępnia
protokołu awaitera wymaganego przez \texttt{co\_await}.
\texttt{future.get()} jest operacją potencjalnie blokującą
bieżący wątek. Istnieją biblioteki dostarczające awaitery dla
przyszłych wyników, ale sam adapter nie zamienia blokującej
operacji w nieblokującą. Naiwne uruchamianie osobnego wątku
tylko po to, by wykonał \texttt{future.wait()} i wznowił korutynę,
może zniweczyć zalety oszczędnego używania wątków.
Dla rzeczywiście nieblokującego oczekiwania na I/O potrzeba
integracji z właściwym mechanizmem asynchronicznym.

\begin{center}
\small
\begin{tabularx}{\textwidth}{@{}lXX@{}}
\toprule
\textbf{Cecha} & \textbf{\texttt{std::async}} & \textbf{Korutyny C++20} \\
\midrule
Co opisuje? & Zadanie ze stanem wyniku & Wykonanie z punktami zawieszenia \\
Tworzy nowy wątek? & \texttt{launch::async}: osobny wątek & Samo \texttt{co\_await}: nie \\
Kiedy zaczyna pracę? & Zależnie od polityki & Zależnie od \texttt{initial\_suspend} i schedulera \\
Odbieranie wyniku & \texttt{future::get()} (może blokować) & Zależy od typu \texttt{Task} i awaitera \\
Wyjątki & Przechowywane w stanie \texttt{future} & Obsługa przez \texttt{unhandled\_exception} i typ wyniku \\
Anulowanie & Brak automatycznego anulowania przez \texttt{future} & Możliwe kooperacyjnie; wymaga protokołu \\
Pamięć/zasoby & Stan asynchroniczny, przy \texttt{async} również wątek & Ramka korutyny, ewentualnie executor \\
Gotowe narzędzie w C++20? & Tak & Składnia tak; ogólny \texttt{Task<T>} nie \\
\bottomrule
\end{tabularx}
\end{center}

\subsection{Ten sam rachunek zapisany na dwa sposoby}

Porównanie dwóch krótkich funkcji pokazuje różnicę semantyczną.

\begin{lstlisting}[caption={std::async uruchamia zadanie w odrębnym wątku (C++11); kompletny program.}]
#include <cassert>
#include <future>

int square(int x) { return x * x; }

int main() {
    auto f = std::async(std::launch::async, square, 7);
    assert(f.get() == 49);
}
\end{lstlisting}

\begin{lstlisting}[caption={Korutyna bez awaitera zewnętrznego nie tworzy wątku; fragment.}]
// Przyjmujemy implementacje Task<T> z poprzedniej sekcji.
Task<int> square_coroutine(int x) {
    co_return x * x;
}

// Task<int> t = square_coroutine(7);
// int result = t.get(); // 49, w watku wywolujacym get()
\end{lstlisting}

W obu przypadkach można otrzymać tę samą wartość, ale nie jest to
to samo zachowanie wykonawcze. \texttt{std::async} z \texttt{launch::async}
zarządza wykonaniem w innym wątku, a pokazany \texttt{Task}
jedynie odkłada wywołanie do momentu ręcznego wznowienia.

\subsection{Pułapka destruktora \texttt{future}}

Ważna różnica dotyczy czasu życia. Gdy ostatni uchwyt do stanu
asynchronicznego utworzonego przez \texttt{std::async} z polityką
\texttt{launch::async} jest niszczony i praca wciąż trwa, destruktor
czeka na zakończenie powiązanego wątku.
Pierwsze z poniższych wywołań zakończy się, zanim rozpocznie się drugie:

\begin{lstlisting}[caption={Zignorowany future wymusza oczekiwanie; fragment.}]
// Funkcje slow_a i slow_b wykonuja dlugie obliczenia.
// std::async(std::launch::async, slow_a);
// std::async(std::launch::async, slow_b);
\end{lstlisting}

Jeżeli chcemy umożliwić nakładanie się zadań, trzeba przechować ich
uchwyty przed oczekiwaniem:
\begin{lstlisting}[caption={Dwa przyszłe wyniki powinny żyć jednocześnie; fragment.}]
// auto fa = std::async(std::launch::async, slow_a);
// auto fb = std::async(std::launch::async, slow_b);
// auto a = fa.get();
// auto b = fb.get();
\end{lstlisting}

Analogicznie w systemie korutyn należy znać politykę niszczenia
\texttt{Task}: usunięcie właściciela zawieszonej korutyny bez
koordynacji z executorem może zostawić wiszący uchwyt. Różne
biblioteki oferują odmienne kontrakty anulowania i oczekiwania.

\subsection{Asynchroniczne I/O a zadania CPU-bound}

Przy \emph{operacji I/O} większość czasu może upływać na oczekiwaniu
na dysk, sieć lub timer. Duża liczba korutyn może opisywać wiele
rozpoczętych operacji, podczas gdy mała pula wątków obsługuje
zdarzenia gotowości i wznowienia. Korzyść nie bierze się z tego,
że \texttt{co\_await} ,,jest szybsze'', lecz z możliwości
nieblokowania wątków podczas oczekiwania.

Przy \emph{operacjach CPU-bound} (obliczenia numeryczne, kompresja,
przetwarzanie obrazów) samo zawieszenie korutyny nie zwiększa
liczby rdzeni wykonujących kod. Tu potrzebne jest rzeczywiste
rozłożenie pracy na wiele wątków (np. \texttt{std::async}, pula
wątków lub OpenMP). Korutyna może być wygodnym \emph{interfejsem}
do takiego executora, ale nie zastępuje go.

\section{Rzeczywista asynchroniczność: korutyna i własny executor}
\label{sec:cpp15-executor}

Żeby zestawienie z \texttt{std::async} nie pozostało czysto
teoretyczne, utworzymy prosty executor z jednym wątkiem roboczym.
\texttt{co\_await executor.schedule()} odłoży kontynuację do
kolejki. Wznowienie wykona \emph{wątek executora}, nie wywołujący.
Wartością zwracaną będzie nasza klasa \texttt{FutureTask<T>},
wykorzystująca wewnętrznie \texttt{std::promise<T>} i
\texttt{std::future<T>}.

W przeciwieństwie do poprzedniego przykładu ramka ma strategię
\emph{samozniszczenia} przez \texttt{final\_suspend = suspend\_never}.
Nie przechowujemy już zewnętrznego uchwytu właścicielskiego:
po zakończeniu korutyna sama domyka swój cykl życia, a klient posiada
wyłącznie \texttt{future}. W tej uproszczonej konstrukcji executor
musi żyć aż do zakończenia zadań, a pojedyncza korutyna jest
wznawiana w kolejce co najwyżej raz.

\begin{lstlisting}[caption={Korutyna przechodząca do innego wątku poprzez executor (C++20); kompletny program.}]
#include <cassert>
#include <condition_variable>
#include <coroutine>
#include <deque>
#include <exception>
#include <future>
#include <mutex>
#include <stdexcept>
#include <thread>
#include <utility>

class Executor {
    std::mutex m_;
    std::condition_variable cv_;
    std::deque<std::coroutine_handle<>> queue_;
    bool closed_ = false;
    std::thread worker_;

    void run() {
        for (;;) {
            std::coroutine_handle<> h;
            {
                std::unique_lock lock(m_);
                cv_.wait(lock, [&] { return closed_ || !queue_.empty(); });
                if (queue_.empty() && closed_) return;
                h = queue_.front();
                queue_.pop_front();
            }
            h.resume(); // H moze sie tu samozniszczyc.
            // Nie korzystamy juz z h po resume().
        }
    }

public:
    Executor() : worker_([this] { run(); }) {}
    Executor(const Executor&) = delete;
    Executor& operator=(const Executor&) = delete;

    ~Executor() {
        {
            std::lock_guard lock(m_);
            closed_ = true;
        }
        cv_.notify_one();
        worker_.join();
    }

    void post(std::coroutine_handle<> h) {
        {
            std::lock_guard lock(m_);
            if (closed_) throw std::runtime_error("Executor closed");
            queue_.push_back(h);
        }
        cv_.notify_one();
    }

    struct Schedule {
        Executor& ex;
        bool await_ready() const noexcept { return false; }
        void await_suspend(std::coroutine_handle<> h) const {
            ex.post(h); // Po przekazaniu h nie uzywamy juz awaitera.
        }
        void await_resume() const noexcept {}
    };

    Schedule schedule() noexcept { return Schedule{*this}; }
};

template<class T>
class FutureTask {
    std::future<T> future_;
public:
    explicit FutureTask(std::future<T> f) : future_(std::move(f)) {}
    FutureTask(FutureTask&&) = default;
    FutureTask(const FutureTask&) = delete;

    T get() { return future_.get(); }

    struct promise_type {
        std::promise<T> result;
        FutureTask get_return_object() {
            return FutureTask{result.get_future()};
        }
        std::suspend_never initial_suspend() noexcept { return {}; }
        std::suspend_never final_suspend() noexcept { return {}; }
        void return_value(T x) { result.set_value(std::move(x)); }
        void unhandled_exception() noexcept {
            try { result.set_exception(std::current_exception()); }
            catch (...) { std::terminate(); }
        }
    };
};

FutureTask<int> square_on(Executor& ex, int x) {
    co_await ex.schedule(); // Tu oddajemy sterowanie.
    co_return x * x;        // Wykona sie w watku roboczym.
}

int main() {
    Executor ex;
    auto a = square_on(ex, 6);
    auto b = square_on(ex, 7);
    assert(a.get() + b.get() == 85);
    // Przed zniszczeniem ex oba zadania sie zakonczyly.
}
\end{lstlisting}

\paragraph{Co to porównanie rzeczywiście dowodzi?}
Kod z \texttt{square\_on} jest asynchroniczny względem wywołującego,
ponieważ po \texttt{co\_await} kontynuację przejmuje inny wątek.
Nie oznacza to jeszcze \emph{równoległego wykonywania} obu obliczeń:
executor ma tylko jednego pracownika, więc zadania są przez niego
wykonywane kolejno. Do równoległego przetwarzania dwóch obliczeń
należałoby dodać drugiego pracownika albo wykorzystać istniejącą
pulę wątków.

Pokazany executor jest \emph{szkicem edukacyjnym}. Nie obsługuje
ogólnego anulowania ani ponownego planowania tej samej korutyny,
nie rozwiązuje wszystkich problemów czasu życia i nie stanowi
gotowej biblioteki. W szczególności klient nie ma uchwytu pozwalającego
anulować i bezpiecznie zniszczyć zawieszone zadanie. W systemie
produkcyjnym własność korutyny, zamknięcie executora i czas życia
powiązanych zasobów muszą stanowić jeden, spójny kontrakt.

\section{Obsługa wyjątków, anulowanie i czas życia danych}
\label{sec:cpp15-errors}

\subsection{Wyjątki a \texttt{unhandled\_exception}}

Nieobsłużony wyjątek z ciała korutyny nie może zostać pozostawiony
bez decyzji. Kompilator wywołuje \texttt{promise.unhandled\_exception()}.
W generatorze zapisaliśmy \texttt{std::exception\_ptr}; przy kolejnym
\texttt{next()} błąd zostanie zgłoszony w kodzie klienta. W naszym
\texttt{FutureTask<T>} wyjątek trafia do stanu \texttt{std::future}
i ujawnia się podczas \texttt{get()}. Jest to analogiczne do
propagacji błędów w \texttt{std::async}, ale sam mechanizm wymagał
jawnej implementacji w \texttt{promise\_type}.

Należy też uwzględnić wyjątki zgłaszane przez operacje awaitera
oraz błędy alokacji ramki. Awaria podczas tworzenia ramki może
wystąpić jeszcze przed wykonaniem ciała korutyny. Typ obietnicy
może opcjonalnie udostępniać \texttt{get\_return\_object\_on\_allocation\_failure()}
razem z odpowiednią, niezgłaszającą wyjątku alokacją, ale nie jest to
mechanizm automatycznie dostępny każdemu typowi korutyny.

\subsection{Referencje przechowywane w ramce}

Argumenty referencyjne nie wydłużają czasu życia obiektów.
Jeżeli korutyna przechowuje referencję do zmiennej lokalnej, która
zniknie podczas zawieszenia, jej późniejsze wznowienie może prowadzić
do użycia nieistniejącego już obiektu.

\begin{lstlisting}[caption={Niebezpieczny czas życia referencji w korutynie; fragment koncepcyjny.}]
// Zakladamy lazy Task<int> z poprzedniej sekcji.
Task<int> read_later(const int& ref) {
    co_return ref;
}

// NIEPOPRAWNE:
// auto task = read_later(123); // 123 znika na koncu tej instrukcji
// int value = task.get();     // odczyt przez wiszaca referencje
\end{lstlisting}

Z reguły bezpieczniej przechowywać wewnątrz ramki wartości
przeznaczone do późniejszego użycia (np. obiekt przekazany przez
wartość), a nie referencje do tymczasowych obiektów.
Trzeba też zwracać uwagę na czas życia domknięć lambda:
korutyna będąca lambdą i wykorzystująca przechwycenia może
po zawieszeniu odwoływać się do już zniszczonego obiektu domknięcia.
Sam fakt przeniesienia wynikowej korutyny nie musi temu zapobiegać.

\subsection{Anulowanie jest kooperacyjne}

\texttt{std::stop\_token} (C++20) umożliwia \emph{żądanie}, a nie
wymuszenie zakończenia zadania. Korutyna może sprawdzać token
przed rozpoczęciem kolejnych faz albo korzystać z awaitera,
który potrafi odpowiednio reagować na anulowanie. Nie wynika
z tego jednak, że każde \texttt{co\_await} automatycznie obsługuje
\texttt{stop\_token}.

\begin{lstlisting}[caption={Kooperacyjne anulowanie korutyny: sprawdzenie żądania w punktach kontrolnych; fragment.}]
// Zakladamy Task<int> z poprzedniej sekcji.
#include <stop_token>

Task<int> calculate(std::stop_token token) {
    if (token.stop_requested()) co_return -1;
    // Pierwszy krok obliczen...
    if (token.stop_requested()) co_return -1;
    // Kolejny krok obliczen...
    co_return 42;
}
\end{lstlisting}

W powyższym przykładzie token jest przekazywany przez wartość.
Sam \texttt{stop\_token} nie steruje czasem życia ramki i nie
wznawia korutyny, która czeka na zdarzenie. Do pełnego anulowania
operacji I/O potrzeba jeszcze protokołu z usługą, która wykonuje
rzeczywistą operację, oraz reguł określających \emph{kto} zniszczy
korutynę i \emph{kiedy} wolno to zrobić.

\section{Pułapki synchronizacji i planowania}

\subsection{Równoczesne wznowienie tej samej korutyny}

Nie można traktować \texttt{coroutine\_handle::resume()} jak
operacji atomowej nadającej się do równoczesnego wywoływania przez
wiele wątków. Trzeba zapewnić, że dana zawieszona korutyna ma
jednego logicznego właściciela uprawnienia do wznowienia.
Ponowne wznowienie korutyny już aktywnej lub zakończonej może
prowadzić do niezdefiniowanego zachowania.

\subsection{Utracone powiadomienie i wyścig publikacji uchwytu}

Awaiter może przekazać uchwyt do innego wątku wewnątrz
\texttt{await\_suspend}. Wówczas druga strona może wznowić korutynę
\emph{zanim} pierwotne wywołanie \texttt{await\_suspend} powróci.
Dlatego po opublikowaniu uchwytu kod awaitera nie powinien korzystać
z obiektów, których czas życia mógł się już zakończyć.

W bardziej złożonych awaiterach trzeba również uniknąć zgubienia
zdarzenia kończącego operację między sprawdzeniem gotowości a
zarejestrowaniem oczekiwania. Poprawne rozwiązania stosują np.
mutex, operacje atomowe albo protokół jednorazowego przejęcia
kontynuacji. Samo \texttt{await\_ready()} nie zastępuje synchronizacji.

\subsection{Brak własności \texttt{coroutine\_handle}}

Gdy planista ma w kolejce jedynie \texttt{coroutine\_handle<>},
nie posiada informacji, kto może zniszczyć ramkę. Usunięcie
obiektu \texttt{Task} podczas oczekiwania może prowadzić do
późniejszego wznowienia już zwolnionej pamięci.
Najbezpieczniejszy projekt definiuje własność od początku:
\emph{Task posiada ramkę}, \emph{executor posiada ramkę} albo
\emph{ramka posiada samą siebie}, z dokładnym protokołem przekazania
własności i zakończenia.

\section{Koszty, optymalizacje i pomiary}
\label{sec:cpp15-performance}

Wydajność należy rozpatrywać w całym systemie, a nie tylko przez
porównanie czasu jednego \texttt{co\_await} i jednego \texttt{std::async}.
\begin{itemize}
\item \textbf{Utworzenie wątku.} \texttt{std::async(launch::async, ...)}
      uruchamia pracę w odrębnym wątku; bardzo duża liczba krótkich
      zadań może generować istotny narzut zarządzania wątkami.
\item \textbf{Ramka.} Korutyna może wymagać alokacji pamięci,
      zachowywania stanu i wznawiania. Kompilator może czasem
      wyeliminować alokację ramki, ale nie jest to gwarancja.
\item \textbf{Executor.} Kolejka, mutexy, przełączanie między wątkami
      oraz współdzielenie pamięci również kosztują.
\item \textbf{Lokalność pamięci.} Duża liczba korutyn z rozproszonymi
      ramkami może prowadzić do dodatkowych nietrafień cache.
\item \textbf{Skalowalność.} Dla I/O ważniejsza może być liczba
      oczekujących operacji niż przepustowość pojedynczego zadania.
      Dla CPU-bound kluczowa jest liczba rdzeni, wielkość porcji pracy
      i ograniczenie dostępu do współdzielonych zasobów.
\end{itemize}

Rzetelny benchmark powinien mierzyć co najmniej: całkowitą
przepustowość, opóźnienia (również wartości odstające), liczbę wątków,
liczbę alokacji oraz maksymalne zużycie pamięci. Porównywanie
\texttt{std::async} z \emph{nieuruchamiającą nowych wątków} korutyną
wykonującą natychmiast prostą arytmetykę nie jest miarodajnym
testem wydajności asynchronicznej.

\subsection{Praktyczna decyzja: kiedy używać którego narzędzia?}

\begin{center}
\small
\begin{tabularx}{\textwidth}{@{}lX@{}}
\toprule
\textbf{Problem} & \textbf{Rozsądny pierwszy wybór} \\
\midrule
Jedno niezależne zadanie CPU z wynikiem & \texttt{std::async(std::launch::async, ...)} \\
Wiele krótkich zadań CPU & Pula wątków, executor lub OpenMP \\
Wiele połączeń I/O z długim oczekiwaniem & Korutyny + asynchroniczny runtime I/O \\
Leniwe generowanie sekwencji & Generator korutynowy \\
Wspólny licznik lub krótka sekcja krytyczna & \texttt{atomic} albo mutex, nie korutyna \\
Anulowalne zadanie działające długo & Przemyślany protokół \texttt{stop\_token} \\
Interfejs biblioteki o złożonych zależnościach & \texttt{Task<T>} nad zdefiniowanym executorem \\
\bottomrule
\end{tabularx}
\end{center}

\section{Związek z OpenMP, \texttt{thread} i \texttt{future}}

Rozszerzony rozdział o współbieżności omawiał \texttt{std::thread},
\texttt{std::jthread}, atomiki, blokady, \texttt{std::async} oraz OpenMP.
Korutyny nie unieważniają tych narzędzi; stanowią kolejną warstwę.
\begin{itemize}
\item \texttt{std::thread} opisuje konkretny wątek wykonania,
      a \texttt{std::jthread} dodatkowo pomaga zarządzać jego
      czasem życia i kooperacyjnym zatrzymywaniem.
\item \texttt{std::async} tworzy stan przyszłego wyniku i
      uruchamia wywołanie zgodnie z przyjętą polityką.
\item OpenMP rozdziela pętle i zadania między wątki obsługiwane
      przez środowisko wykonawcze OpenMP; szczególnie przydaje się
      w obliczeniach CPU-bound i kodzie numerycznym.
\item Korutyny pozwalają wygodnie opisać wieloetapową logikę
      oczekiwania i kontynuacji. Ich wydajność i stopień równoległości
      zależą od dostarczonego środowiska wykonawczego.
\end{itemize}

Możliwe jest użycie \emph{korutyn nad pulą wątków}: \texttt{co\_await}
przekazuje kontynuację do kolejki, a pracownik puli później ją wznawia.
Nie należy jednak od razu łączyć wszystkich mechanizmów.
Najpierw określamy, czy program czeka głównie na I/O, czy wykonuje
obliczenia, a następnie wybieramy model wykonania i własności zasobów.

\section{Najczęstsze błędy projektowe}

\begin{enumerate}
\item Zakładanie, że \texttt{co\_await} uruchamia nowy wątek.
\item Próba zwrócenia bezpośrednio \texttt{std::future<T>}
      z własnej korutyny bez odpowiedniego \texttt{promise\_type}.
\item Dwukrotne wywołanie \texttt{destroy()} dla tej samej ramki.
\item Wywołanie \texttt{resume()} po osiągnięciu
      \texttt{final\_suspend}.
\item Pozostawienie referencji lub przechwycenia lambdy do obiektu,
      który znika podczas zawieszenia.
\item Zniszczenie \texttt{Task} mimo tego, że executor posiada
      w kolejce uchwyt do jego ramki.
\item Mylenie automatycznej propagacji wyjątków w \texttt{std::async}
      z obowiązkami implementatora \texttt{promise\_type}.
\item Traktowanie \texttt{std::stop\_token} jako mechanizmu
      przymusowego przerwania dowolnego \texttt{co\_await}.
\item Równoczesne wznawianie tej samej korutyny z kilku wątków.
\item Uznanie, że zawieszenie i wznowienie zawsze kosztuje mniej
      niż zwykłe wywołanie funkcji.
\item Uruchamianie \texttt{std::async} dla tysięcy mikro-zadań
      bez rozważenia puli wątków.
\item Zwracanie z \texttt{await\_suspend} po opublikowaniu uchwytu
      i dalsze używanie potencjalnie zniszczonego awaitera.
\end{enumerate}

\section*{Podsumowanie rozdziału}
\addcontentsline{toc}{section}{Podsumowanie rozdziału}

Korutyny C++20 są mechanizmem \emph{zawieszenia i kontynuacji},
a nie samodzielnym środowiskiem asynchronicznym. Generator używa
\texttt{co\_yield}, aby leniwie udostępniać wartości. Typ zadania
używa \texttt{co\_return}, obiektu \texttt{promise\_type} i awaiterów,
aby opisać wykonanie z wynikiem i wyjątkami. \texttt{co\_await}
przekazuje kontrolę zgodnie z protokołem awaitera, a dopiero executor
lub inne środowisko może wznowić kod w osobnym wątku.

\texttt{std::async} jest prostym narzędziem do uruchomienia funkcji
i odebrania przyszłego wyniku. Korutyny są lepszym budulcem dla
złożonych przepływów wieloetapowych, szczególnie I/O, ale wymagają
większej odpowiedzialności za pamięć, synchronizację, anulowanie
i czas życia ramki. W obu przypadkach należy rozdzielać trzy pytania:
\emph{kto wykonuje kod, kto posiada stan i kto decyduje o zakończeniu}.

\section*{Literatura i dokumentacja}
\addcontentsline{toc}{section}{Literatura i dokumentacja}
\begin{itemize}
\item \url{https://en.cppreference.com/w/cpp/language/coroutines} --- model korutyn C++20.
\item \url{https://en.cppreference.com/w/cpp/coroutine/coroutine_handle} --- uchwyty i zarządzanie ramką.
\item \url{https://en.cppreference.com/w/cpp/coroutine/suspend_always} --- podstawowe awaitery.
\item \url{https://en.cppreference.com/w/cpp/thread/async} --- \texttt{std::async} i polityki uruchamiania.
\item \url{https://en.cppreference.com/w/cpp/thread/future} --- \texttt{std::future}.
\item \url{https://en.cppreference.com/w/cpp/coroutine/generator} --- \texttt{std::generator} (C++23).
\item Lewis Baker, \emph{C++ Coroutines: Understanding the Compiler Transformations}
      --- seria publikacji o modelu wykonania korutyn.
\item David Mazières, \emph{My Tutorial and Take on C++20 Coroutines}, 2021.
\end{itemize}

% Listing oznaczony jako „kompletny program” mozna skompilowac oddzielnie.

\chapter{Korutyny, zadania asynchroniczne i \texttt{std::async}}
\label{chap:cpp15-coroutines}

\section{Po co C++ potrzebuje korutyn?}

Zwykła funkcja ma zasadniczo jeden punkt wejścia. Może wywoływać inne
funkcje, ale po powrocie oddaje sterowanie wywołującemu i nie posiada
mechanizmu późniejszego wznowienia dokładnie w miejscu, w którym
zakończyła wykonywanie. \emph{Korutyna} (ang. \emph{coroutine}) jest
funkcją, której wykonanie można \emph{zawiesić}, a następnie
\emph{wznowić}, zachowując wybrane zmienne lokalne i miejsce kontynuacji.

Korutyny C++20 wspierają między innymi:
\begin{itemize}
\item \emph{generatory} --- wytwarzanie kolejnych wartości na żądanie,
      bez budowania całej kolekcji;
\item \emph{zadania asynchroniczne} --- opis operacji, których kolejne
      etapy uruchamiane są po zakończeniu wcześniejszych;
\item \emph{kooperacyjne planowanie} --- oddawanie sterowania
      harmonogramowi zadań;
\item \emph{maszyny stanów} --- zachowywanie stanu pomiędzy wznowieniami
      bez ręcznego kodowania każdej możliwej kontynuacji.
\end{itemize}

Najważniejsza zasada brzmi:
\begin{center}
\emph{Korutyna nie jest wątkiem, a \texttt{co\_await} nie oznacza
automatycznie ,,uruchom w nowym wątku''.}
\end{center}

Korutyna określa \emph{jak zatrzymać i kontynuować obliczenie};
\texttt{std::async}, \texttt{std::thread} i pula wątków dotyczą również
\emph{gdzie i kiedy} wykonywać pracę. W szczególności kod korutyny
może działać od początku do końca w jednym wątku, nawet jeśli
wielokrotnie się zawiesza. Odwrotnie, korutyna może być wznawiana
przez różne wątki, jeżeli tak zdecyduje dostarczony mechanizm planowania.

\subsection{Trzy słowa kluczowe}

\begin{center}\small
\begin{tabularx}{\textwidth}{@{}p{0.32\textwidth}X@{}}
\toprule
Konstrukcja & Działanie w korutynie \\
\midrule
\code{co_await expr} & Przekazuje sterowanie obiektowi oczekującemu; może zawiesić korutynę bez blokowania wątku. \\
\code{co_yield value} & Udostępnia kolejną wartość i zwykle zawiesza generator do wznowienia. \\
\code{co_return value} & Kończy korutynę i przekazuje wynik zgodnie z protokołem jej typu zwracanego. \\
\code{await_ready()} & Określa, czy wynik jest już gotowy i można pominąć zawieszenie. \\
\code{await_suspend(h)} & Obsługuje zawieszenie i uchwyt \code{h}; może zaplanować późniejsze wznowienie. \\
\code{await_resume()} & Zwraca rezultat oczekiwania albo zgłasza zapisany wyjątek po wznowieniu. \\
\bottomrule
\end{tabularx}
\end{center}

C++20 wprowadził następujące konstrukcje składniowe:
\begin{description}
\item[\texttt{co\_await expr}] Zastosuj protokół oczekiwania do wyrażenia;
      zależnie od niego możesz zawiesić korutynę i później ją wznowić.
\item[\texttt{co\_yield expr}] Udostępnij kolejną wartość generatora;
      składniowo jest to szczególna postać oczekiwania, wykorzystująca
      metodę \texttt{yield\_value} obiektu \texttt{promise\_type}.
\item[\texttt{co\_return expr}] Zakończ korutynę i przekaż wynik
      przez odpowiednie \texttt{return\_value} albo \texttt{return\_void}.
\end{description}

Funkcja zawierająca jeden z tych operatorów (w odpowiednim kontekście)
rozpoznawana jest jako korutyna. Typ jej zwracanej wartości musi
udostępniać właściwy \emph{typ obietnicy} \texttt{promise\_type}
(lub odpowiednią specjalizację \texttt{std::coroutine\_traits}).
To wymóg mechanizmu korutyn, a nie znanego już \texttt{std::promise}.

\begin{lstlisting}[caption={Niepoprawna intuicja: std::future nie jest automatycznie typem korutyny (C++20); fragment celowo błędny.}]
#include <future>

// Standardowe std::future<int> nie zapewnia promise_type:
// std::future<int> calculate() { co_return 42; } // BLAD!
\end{lstlisting}

W C++20 biblioteka standardowa dostarcza \emph{mechanizm} korutyn
(nagłówek \texttt{<coroutine>}), ale nie dostarcza gotowego, ogólnego
\texttt{Task<T>} ani standardowego executora. Takie typy definiuje
biblioteka lub programista. \texttt{std::generator} jest rozwiązaniem
późniejszym, wprowadzonym dopiero w C++23, a jego wsparcie zależy
od użytej biblioteki standardowej.

\section{Różnica między zawieszeniem a zablokowaniem wątku}
\label{sec:cpp15-suspend-block}

Rozważmy dwa sposoby wstrzymania działania programu:
\begin{itemize}
\item \texttt{std::this\_thread::sleep\_for(100ms)} \emph{blokuje
      bieżący wątek} na pewien czas; sam wątek pozostaje zajętym
      zasobem systemowym, choć może nie zużywać wówczas czasu CPU;
\item dobrze zaprojektowane \texttt{co\_await} \emph{zawiesza
      korutynę}, oddając wątek do wykonywania innych zadań;
      wznowienie może nastąpić po nadejściu wyniku operacji I/O,
      komunikatu lub upływie czasu.
\end{itemize}

Nie każde \texttt{co\_await} daje zawieszenie: awaiter może zgłosić,
że wynik jest już gotowy. Również samo użycie \texttt{co\_await}
nie gwarantuje nieblokującego I/O --- zależy to od jego implementacji.

\begin{center}
\small
\begin{tabularx}{\textwidth}{@{}lXX@{}}
\toprule
\textbf{Właściwość} & \textbf{Blokada wątku} & \textbf{Zawieszenie korutyny} \\
\midrule
Zachowanie stanu & Na stosie wywołań wątku & Zwykle w ramce korutyny \\
Wątek potrzebny podczas oczekiwania & Tak, istnieje i czeka & Nie jest konieczny \emph{dla samego oczekiwania} \\
Wznowienie & Planista systemu/odblokowanie & Mechanizm awaitera i właściciel kontynuacji \\
Koszt & M.in. wątek i jego stos & M.in. ramka i obsługa zawieszenia \\
\bottomrule
\end{tabularx}
\end{center}

\section{Ramka korutyny i jej cykl życia}
\label{sec:cpp15-frame}

Kompilator przekształca korutynę w obiekt reprezentujący jej stan.
Ten stan nazywamy \emph{ramką korutyny} (ang. \emph{coroutine frame}).
Może ona obejmować obiekt \texttt{promise\_type}, parametry i stan
potrzebny do kontynuowania obliczenia po zawieszeniu. To, które
zmienne lokalne trzeba zachować, wynika z przebiegu programu.

W ogólnym modelu wywołanie korutyny przebiega w następujących etapach:
\begin{enumerate}
\item Alokowana jest ramka, chyba że implementacja legalnie
      wyeliminuje tę alokację (np. przez zagnieżdżenie ramki w innym obiekcie).
\item Kopiowane lub przenoszone są parametry odpowiednio do ich rodzajów;
      \emph{parametry referencyjne pozostają referencjami}.
\item Konstruowany jest obiekt obietnicy \texttt{promise\_type}.
\item Wywoływana jest metoda \texttt{get\_return\_object()}.
\item Wykonywany jest \texttt{co\_await promise.initial\_suspend()}.
\item Uruchamiane jest ciało korutyny (od razu lub po wznowieniu).
\item \texttt{co\_return} uruchamia \texttt{return\_value} lub
      \texttt{return\_void}; nieobsłużony wyjątek trafia do
      \texttt{unhandled\_exception()}.
\item Niszczone są obiekty lokalne, a korutyna wykonuje końcowy
      \texttt{co\_await} na wyniku \texttt{promise.final\_suspend()}.
\item Ramka zostaje ostatecznie zniszczona we właściwym momencie.
\end{enumerate}

Kolejność i czas \emph{zwolnienia ramki} zależą od przyjętej polityki
własności. Dla \texttt{final\_suspend = suspend\_always} ramkę zwykle
niszczy właściciel uchwytu, nawet po zakończeniu ciała korutyny.
Przy \texttt{final\_suspend = suspend\_never} mechanizm korutyny
może zakończyć się bez pozostawiania ramki zawieszonej do późniejszego
zniszczenia przez właściciela. Tych strategii nie wolno mieszać.

\subsection{Początek i koniec: \texttt{initial\_suspend} i \texttt{final\_suspend}}

Dwa najczęściej używane typy z nagłówka \texttt{<coroutine>} to:
\begin{description}
\item[\texttt{std::suspend\_never}] Oczekiwanie jest natychmiast
      gotowe i nie powoduje zawieszenia.
\item[\texttt{std::suspend\_always}] Oczekiwanie zawsze prowadzi
      do zawieszenia, dopóki ktoś nie wznowi korutyny.
\end{description}

Jeżeli \texttt{initial\_suspend()} zwraca \texttt{suspend\_always},
korutyna jest \emph{lazy}: nie rozpoczyna wykonywania ciała przy samym
wywołaniu funkcji. Natomiast \texttt{suspend\_never} oznacza
\emph{eager start}: ciało może rozpocząć się natychmiast w wątku
wywołującym i pracować aż do pierwszego rzeczywistego zawieszenia.

\textbf{Ważne:} nie wolno wywoływać \texttt{resume()} na korutynie,
która jest zawieszona w swoim \texttt{final\_suspend}. Jest to
niezdefiniowane zachowanie. Uchwyt do takiej ramki może nadal służyć
do \texttt{destroy()}, jeżeli ramka nie została już zniszczona.

\section{Uchwyt \texttt{std::coroutine\_handle} i własność ramki}

\texttt{std::coroutine\_handle<Promise>} jest lekkim uchwytem
umożliwiającym m.in. wznowienie korutyny i dostęp do jej obietnicy.
Sam \texttt{coroutine\_handle} \emph{nie zarządza automatycznie
czasem życia ramki}. Skopiowanie uchwytu nie jest kopią korutyny,
a zniszczenie uchwytu nie wywołuje automatycznie \texttt{destroy()}.
Dlatego wygodna klasa będąca właścicielem ramki powinna stosować RAII,
uniemożliwiać kopiowanie oraz poprawnie obsługiwać przenoszenie.

\begin{lstlisting}[caption={Operacje na uchwycie; fragment ilustracyjny.}]
#include <coroutine>

using Handle = std::coroutine_handle<>;

// Zakladamy, ze h wskazuje na istniejaca, zawieszona korutyne.
// if (h && !h.done()) h.resume();
// h.destroy(); // dokladnie raz, przez wlasciciela ramki
// Po destroy() wszystkie kopie h sa wiszacymi uchwytami.
\end{lstlisting}

Uchwyt może zostać przekazany do innego wątku lub executora, ale
wówczas trzeba zagwarantować, że ramka nie zostanie zniszczona,
zanim planista faktycznie wznowi zawieszoną korutynę. Jest to jedna
z głównych trudności implementacyjnych asynchronicznego API C++20.

\section{Własny generator: \texttt{co\_yield} od podstaw}
\label{sec:cpp15-generator}

Zbudujemy teraz typ \texttt{Generator}, który produkuje liczby całkowite
kolejno, \emph{na żądanie}. Poniższy przykład jest samowystarczalny.
Prywatny uchwyt do ramki stanowi jedynego właściciela zasobu;
\texttt{next()} wznawia korutynę tylko wtedy, gdy nie zakończyła
wykonywania. Własną obsługą wyjątku uzupełniamy przykład o
poprawną propagację błędów.

\begin{lstlisting}[caption={Generator liczb całkowitych: co\_yield i RAII (C++20); kompletny program.}]
#include <cassert>
#include <coroutine>
#include <exception>
#include <limits>
#include <stdexcept>
#include <utility>

class Generator {
public:
    struct promise_type {
        int current = 0;
        std::exception_ptr error;

        Generator get_return_object() {
            return Generator{Handle::from_promise(*this)};
        }
        std::suspend_always initial_suspend() noexcept { return {}; }
        std::suspend_always final_suspend() noexcept { return {}; }
        std::suspend_always yield_value(int x) noexcept {
            current = x;
            return {};
        }
        void return_void() noexcept {}
        void unhandled_exception() noexcept {
            error = std::current_exception();
        }
    };

    using Handle = std::coroutine_handle<promise_type>;

    explicit Generator(Handle h) noexcept : handle_(h) {}
    Generator(const Generator&) = delete;
    Generator& operator=(const Generator&) = delete;

    Generator(Generator&& other) noexcept
        : handle_(std::exchange(other.handle_, {})),
          ready_(std::exchange(other.ready_, false)) {}

    Generator& operator=(Generator&& other) noexcept {
        if (this != &other) {
            if (handle_) handle_.destroy();
            handle_ = std::exchange(other.handle_, {});
            ready_ = std::exchange(other.ready_, false);
        }
        return *this;
    }

    ~Generator() {
        if (handle_) handle_.destroy();
    }

    bool next() {
        ready_ = false;
        if (!handle_ || handle_.done()) return false;
        handle_.resume();
        if (handle_.promise().error)
            std::rethrow_exception(handle_.promise().error);
        ready_ = !handle_.done();
        return ready_;
    }

    int value() const {
        if (!ready_)
            throw std::logic_error("No current value");
        return handle_.promise().current;
    }

private:
    Handle handle_{};
    bool ready_ = false;
};

Generator squares(int n) {
    for (int i = 1; i <= n; ++i) {
        if (i > std::numeric_limits<int>::max() / i)
            throw std::overflow_error("square outside int");
        co_yield i * i;
    }
}

int main() {
    auto gen = squares(4); // Cialo jeszcze sie nie wykonalo.
    int sum = 0;
    while (gen.next()) sum += gen.value();
    assert(sum == 1 + 4 + 9 + 16);
    assert(!gen.next());
}
\end{lstlisting}

\subsection{Co dzieje się podczas \texttt{co\_yield}?}

Instrukcja \texttt{co\_yield i * i} powoduje wywołanie
\texttt{promise.yield\_value(i * i)} i zastosowanie mechanizmu
\texttt{co\_await} do otrzymanego awaitera. W naszym przypadku
\texttt{yield\_value} zapisuje wartość w obietnicy i zwraca
\texttt{suspend\_always}. Funkcja \texttt{next()} otrzymuje wtedy
sterowanie z powrotem i udostępnia zapisany element. Następne
\texttt{next()} kontynuuje pętlę od poprzedniego miejsca zawieszenia.

Generator posiada tylko jeden bieżący wynik. Nie przechowuje
całego ciągu jak \texttt{std::vector<int>}. To użyteczne przy dużych
lub potencjalnie nieskończonych sekwencjach. \emph{Nie oznacza to}
jednak automatycznie mniejszego kosztu obliczeń --- dochodzą koszty
ramki i kolejnych wznowień.

\nota{Przedstawiony \texttt{Generator} jest celowo minimalny: nie
udostępnia interfejsu iteratorów \texttt{begin()/end()}, nie jest
bezpieczny do jednoczesnego wznawiania z kilku wątków i nie powinien
być stosowany jako gotowa biblioteka produkcyjna.}

\section{Protokół \texttt{co\_await}: \texttt{await\_ready},
\texttt{await\_suspend}, \texttt{await\_resume}}
\label{sec:cpp15-awaiter}

Operator \texttt{co\_await} nie jest synonimem oczekiwania na
\texttt{future}. Kompilator uzyskuje obiekt \emph{awaiter}, którego
interfejs ma trzy fundamentalne operacje:
\begin{description}
\item[\texttt{await\_ready()}] Zwraca wartość logiczną: jeśli
      \texttt{true}, zawieszenie jest pomijane.
\item[\texttt{await\_suspend(handle)}] Jeżeli \texttt{await\_ready()}
      było fałszywe, metoda podejmuje decyzję o kontynuacji.
      Dopuszczalne postacie zwracają \texttt{void}, \texttt{bool}
      albo inny \texttt{std::coroutine\_handle}.
\item[\texttt{await\_resume()}] Wyznacza wynik całego wyrażenia
      \texttt{co\_await}; może też zgłosić wyjątek.
\end{description}

Gdy \texttt{await\_suspend} zwraca \texttt{bool}, wartość
\texttt{false} oznacza kontynuację bez pozostawania w zawieszeniu,
a \texttt{true} --- pozostawienie korutyny w stanie zawieszonym.
Przy zwrocie \texttt{void} korutyna pozostaje zawieszona do późniejszego
wznowienia. Zwrócenie innego uchwytu powoduje \emph{symmetric transfer}:
sterowanie przechodzi do wskazanej korutyny.

Nie każdy typ bezpośrednio dostarcza trzy wymienione metody: język
obsługuje też przeciążenie \texttt{operator co\_await} oraz
\texttt{await\_transform} zdefiniowane w \texttt{promise\_type}.
Sposób uzyskania awaitera zależy od tych reguł.

\subsection{Czysto lokalne zawieszenie bez nowego wątku}

Poniższy typ liczy wznowienia. Wykorzystamy go z prostą korutyną,
której właściciel ręcznie wywołuje \texttt{resume()}.

\begin{lstlisting}[caption={co\_await z własnym awaiterem: żadnego nowego wątku (C++20); kompletny program.}]
#include <cassert>
#include <coroutine>
#include <exception>
#include <utility>

struct Stepper {
    struct promise_type {
        Stepper get_return_object() {
            return Stepper{handle::from_promise(*this)};
        }
        std::suspend_always initial_suspend() noexcept { return {}; }
        std::suspend_always final_suspend() noexcept { return {}; }
        void return_void() noexcept {}
        void unhandled_exception() { std::terminate(); }
    };

    using handle = std::coroutine_handle<promise_type>;
    handle h;
    explicit Stepper(handle x) : h(x) {}
    Stepper(const Stepper&) = delete;
    Stepper(Stepper&& other) noexcept
        : h(std::exchange(other.h, {})) {}
    ~Stepper() { if (h) h.destroy(); }
    bool resume() {
        if (!h || h.done()) return false;
        h.resume();
        return !h.done();
    }
};

struct Pause {
    int& count;
    bool await_ready() const noexcept { return false; }
    void await_suspend(std::coroutine_handle<>) const noexcept {
        ++count; // Zawiesz, az wlasciciel wywola resume().
    }
    void await_resume() const noexcept {}
};

Stepper example(int& count) {
    ++count;                 // wykona sie przy pierwszym resume
    co_await Pause{count};  // pierwsze zawieszenie
    ++count;                 // wykona sie przy drugim resume
}

int main() {
    int count = 0;
    auto step = example(count);
    assert(count == 0);
    assert(step.resume() && count == 2);
    assert(!step.resume() && count == 3);
}
\end{lstlisting}

Kod niczego nie oblicza równolegle. Wszystkie instrukcje wykonał
wątek \texttt{main}. Korutyna jedynie oddała sterowanie, zachowując
miejsce wznowienia.

\section{Własne \texttt{Task<T>}: leniwe obliczenie a asynchroniczność}
\label{sec:cpp15-task}

Zaprojektujemy najmniejszy rozsądny typ \texttt{Task<T>} zwracający
wynik i obsługujący wyjątek. Celowo nie implementujemy jeszcze
harmonogramu: pozwala to zrozumieć, że \emph{sam typ korutyny nie jest
schedulerem}. \texttt{get()} w tym przykładzie działa wyłącznie
wtedy, gdy ciało może zakończyć się po pojedynczym wznowieniu.

\begin{lstlisting}[caption={Lazy Task<T> bez schedulera (C++20); kompletny program.}]
#include <cassert>
#include <coroutine>
#include <exception>
#include <optional>
#include <stdexcept>
#include <utility>

template<class T>
class Task {
public:
    struct promise_type {
        std::optional<T> result;
        std::exception_ptr error;
        Task get_return_object() {
            return Task{Handle::from_promise(*this)};
        }
        std::suspend_always initial_suspend() noexcept { return {}; }
        std::suspend_always final_suspend() noexcept { return {}; }
        void return_value(T x) { result.emplace(std::move(x)); }
        void unhandled_exception() noexcept {
            error = std::current_exception();
        }
    };
    using Handle = std::coroutine_handle<promise_type>;

    explicit Task(Handle h) noexcept : h_(h) {}
    Task(const Task&) = delete;
    Task& operator=(const Task&) = delete;
    Task(Task&& x) noexcept
        : h_(std::exchange(x.h_, {})),
          taken_(std::exchange(x.taken_, false)) {}
    ~Task() { if (h_) h_.destroy(); }

    T get() {
        if (!h_) throw std::logic_error("Empty task");
        if (taken_) throw std::logic_error("Result already consumed");
        if (!h_.done()) h_.resume();
        if (!h_.done())
            throw std::logic_error("Task needs an external scheduler");
        taken_ = true;
        if (h_.promise().error)
            std::rethrow_exception(h_.promise().error);
        return std::move(*h_.promise().result);
    }

private:
    Handle h_{};
    bool taken_ = false;
};

Task<int> calculate(int x) {
    co_return x * x;
}

int main() {
    auto task = calculate(7); // lazy, obliczenie jeszcze nie ruszylo
    assert(task.get() == 49);
}
\end{lstlisting}

Ten \texttt{Task<T>} jest pedagogiczny: nie ma przeciążenia
\texttt{operator co\_await}, nie wspiera wielokrotnego pobierania
wyniku (drugie \texttt{get()} zgłasza błąd) i nie potrafi samodzielnie
czekać na zdarzenia zewnętrzne.
W szczególności \texttt{get()} nie może po prostu wielokrotnie
wywoływać \texttt{resume()} na rzeczywistym zadaniu I/O: wznowienie
przed ukończeniem operacji mogłoby naruszyć protokół awaitera.
Produkcyjne \texttt{Task<T>} wymaga m.in. protokołu kontynuacji,
jasnej semantyki własności i poprawnej synchronizacji z executorem.

\section{Korutyna a \texttt{std::async}: fundamentalne porównanie}
\label{sec:cpp15-async-comparison}

\texttt{std::async} (C++11, nagłówek \texttt{<future>}) jest gotową
funkcją uruchamiającą zadanie i zwracającą \texttt{std::future<T>}.
Nie wymaga od użytkownika implementowania \texttt{promise\_type}.
Przy polityce \texttt{std::launch::async} funkcja pracuje w odrębnym
wątku; polityka \texttt{std::launch::deferred} odracza jej uruchomienie
do pierwszego nieczasowego oczekiwania na wynik.

Z kolei korutyna C++20 definiuje sposób zapisywania \emph{kontynuacji}.
Nawet jeśli wynik ma charakter asynchroniczny, \emph{faktyczne
uruchamianie pracy} zależy od awaitera, event loop, schedulera lub
innego mechanizmu zewnętrznego. Nie istnieje standardowa ogólna
reguła ,,jedna korutyna = jeden wątek''.


\paragraph{Czy można napisać \texttt{co\_await future}?}
Nie wprost: w C++20 standardowy \texttt{std::future<T>} nie udostępnia
protokołu awaitera wymaganego przez \texttt{co\_await}.
\texttt{future.get()} jest operacją potencjalnie blokującą
bieżący wątek. Istnieją biblioteki dostarczające awaitery dla
przyszłych wyników, ale sam adapter nie zamienia blokującej
operacji w nieblokującą. Naiwne uruchamianie osobnego wątku
tylko po to, by wykonał \texttt{future.wait()} i wznowił korutynę,
może zniweczyć zalety oszczędnego używania wątków.
Dla rzeczywiście nieblokującego oczekiwania na I/O potrzeba
integracji z właściwym mechanizmem asynchronicznym.

\begin{center}
\small
\begin{tabularx}{\textwidth}{@{}lXX@{}}
\toprule
\textbf{Cecha} & \textbf{\texttt{std::async}} & \textbf{Korutyny C++20} \\
\midrule
Co opisuje? & Zadanie ze stanem wyniku & Wykonanie z punktami zawieszenia \\
Tworzy nowy wątek? & \texttt{launch::async}: osobny wątek & Samo \texttt{co\_await}: nie \\
Kiedy zaczyna pracę? & Zależnie od polityki & Zależnie od \texttt{initial\_suspend} i schedulera \\
Odbieranie wyniku & \texttt{future::get()} (może blokować) & Zależy od typu \texttt{Task} i awaitera \\
Wyjątki & Przechowywane w stanie \texttt{future} & Obsługa przez \texttt{unhandled\_exception} i typ wyniku \\
Anulowanie & Brak automatycznego anulowania przez \texttt{future} & Możliwe kooperacyjnie; wymaga protokołu \\
Pamięć/zasoby & Stan asynchroniczny, przy \texttt{async} również wątek & Ramka korutyny, ewentualnie executor \\
Gotowe narzędzie w C++20? & Tak & Składnia tak; ogólny \texttt{Task<T>} nie \\
\bottomrule
\end{tabularx}
\end{center}

\subsection{Ten sam rachunek zapisany na dwa sposoby}

Porównanie dwóch krótkich funkcji pokazuje różnicę semantyczną.

\begin{lstlisting}[caption={std::async uruchamia zadanie w odrębnym wątku (C++11); kompletny program.}]
#include <cassert>
#include <future>

int square(int x) { return x * x; }

int main() {
    auto f = std::async(std::launch::async, square, 7);
    assert(f.get() == 49);
}
\end{lstlisting}

\begin{lstlisting}[caption={Korutyna bez awaitera zewnętrznego nie tworzy wątku; fragment.}]
// Przyjmujemy implementacje Task<T> z poprzedniej sekcji.
Task<int> square_coroutine(int x) {
    co_return x * x;
}

// Task<int> t = square_coroutine(7);
// int result = t.get(); // 49, w watku wywolujacym get()
\end{lstlisting}

W obu przypadkach można otrzymać tę samą wartość, ale nie jest to
to samo zachowanie wykonawcze. \texttt{std::async} z \texttt{launch::async}
zarządza wykonaniem w innym wątku, a pokazany \texttt{Task}
jedynie odkłada wywołanie do momentu ręcznego wznowienia.

\subsection{Pułapka destruktora \texttt{future}}

Ważna różnica dotyczy czasu życia. Gdy ostatni uchwyt do stanu
asynchronicznego utworzonego przez \texttt{std::async} z polityką
\texttt{launch::async} jest niszczony i praca wciąż trwa, destruktor
czeka na zakończenie powiązanego wątku.
Pierwsze z poniższych wywołań zakończy się, zanim rozpocznie się drugie:

\begin{lstlisting}[caption={Zignorowany future wymusza oczekiwanie; fragment.}]
// Funkcje slow_a i slow_b wykonuja dlugie obliczenia.
// std::async(std::launch::async, slow_a);
// std::async(std::launch::async, slow_b);
\end{lstlisting}

Jeżeli chcemy umożliwić nakładanie się zadań, trzeba przechować ich
uchwyty przed oczekiwaniem:
\begin{lstlisting}[caption={Dwa przyszłe wyniki powinny żyć jednocześnie; fragment.}]
// auto fa = std::async(std::launch::async, slow_a);
// auto fb = std::async(std::launch::async, slow_b);
// auto a = fa.get();
// auto b = fb.get();
\end{lstlisting}

Analogicznie w systemie korutyn należy znać politykę niszczenia
\texttt{Task}: usunięcie właściciela zawieszonej korutyny bez
koordynacji z executorem może zostawić wiszący uchwyt. Różne
biblioteki oferują odmienne kontrakty anulowania i oczekiwania.

\subsection{Asynchroniczne I/O a zadania CPU-bound}

Przy \emph{operacji I/O} większość czasu może upływać na oczekiwaniu
na dysk, sieć lub timer. Duża liczba korutyn może opisywać wiele
rozpoczętych operacji, podczas gdy mała pula wątków obsługuje
zdarzenia gotowości i wznowienia. Korzyść nie bierze się z tego,
że \texttt{co\_await} ,,jest szybsze'', lecz z możliwości
nieblokowania wątków podczas oczekiwania.

Przy \emph{operacjach CPU-bound} (obliczenia numeryczne, kompresja,
przetwarzanie obrazów) samo zawieszenie korutyny nie zwiększa
liczby rdzeni wykonujących kod. Tu potrzebne jest rzeczywiste
rozłożenie pracy na wiele wątków (np. \texttt{std::async}, pula
wątków lub OpenMP). Korutyna może być wygodnym \emph{interfejsem}
do takiego executora, ale nie zastępuje go.

\section{Rzeczywista asynchroniczność: korutyna i własny executor}
\label{sec:cpp15-executor}

Żeby zestawienie z \texttt{std::async} nie pozostało czysto
teoretyczne, utworzymy prosty executor z jednym wątkiem roboczym.
\texttt{co\_await executor.schedule()} odłoży kontynuację do
kolejki. Wznowienie wykona \emph{wątek executora}, nie wywołujący.
Wartością zwracaną będzie nasza klasa \texttt{FutureTask<T>},
wykorzystująca wewnętrznie \texttt{std::promise<T>} i
\texttt{std::future<T>}.

W przeciwieństwie do poprzedniego przykładu ramka ma strategię
\emph{samozniszczenia} przez \texttt{final\_suspend = suspend\_never}.
Nie przechowujemy już zewnętrznego uchwytu właścicielskiego:
po zakończeniu korutyna sama domyka swój cykl życia, a klient posiada
wyłącznie \texttt{future}. W tej uproszczonej konstrukcji executor
musi żyć aż do zakończenia zadań, a pojedyncza korutyna jest
wznawiana w kolejce co najwyżej raz.

\begin{lstlisting}[caption={Korutyna przechodząca do innego wątku poprzez executor (C++20); kompletny program.}]
#include <cassert>
#include <condition_variable>
#include <coroutine>
#include <deque>
#include <exception>
#include <future>
#include <mutex>
#include <stdexcept>
#include <thread>
#include <utility>

class Executor {
    std::mutex m_;
    std::condition_variable cv_;
    std::deque<std::coroutine_handle<>> queue_;
    bool closed_ = false;
    std::thread worker_;

    void run() {
        for (;;) {
            std::coroutine_handle<> h;
            {
                std::unique_lock lock(m_);
                cv_.wait(lock, [&] { return closed_ || !queue_.empty(); });
                if (queue_.empty() && closed_) return;
                h = queue_.front();
                queue_.pop_front();
            }
            h.resume(); // H moze sie tu samozniszczyc.
            // Nie korzystamy juz z h po resume().
        }
    }

public:
    Executor() : worker_([this] { run(); }) {}
    Executor(const Executor&) = delete;
    Executor& operator=(const Executor&) = delete;

    ~Executor() {
        {
            std::lock_guard lock(m_);
            closed_ = true;
        }
        cv_.notify_one();
        worker_.join();
    }

    void post(std::coroutine_handle<> h) {
        {
            std::lock_guard lock(m_);
            if (closed_) throw std::runtime_error("Executor closed");
            queue_.push_back(h);
        }
        cv_.notify_one();
    }

    struct Schedule {
        Executor& ex;
        bool await_ready() const noexcept { return false; }
        void await_suspend(std::coroutine_handle<> h) const {
            ex.post(h); // Po przekazaniu h nie uzywamy juz awaitera.
        }
        void await_resume() const noexcept {}
    };

    Schedule schedule() noexcept { return Schedule{*this}; }
};

template<class T>
class FutureTask {
    std::future<T> future_;
public:
    explicit FutureTask(std::future<T> f) : future_(std::move(f)) {}
    FutureTask(FutureTask&&) = default;
    FutureTask(const FutureTask&) = delete;

    T get() { return future_.get(); }

    struct promise_type {
        std::promise<T> result;
        FutureTask get_return_object() {
            return FutureTask{result.get_future()};
        }
        std::suspend_never initial_suspend() noexcept { return {}; }
        std::suspend_never final_suspend() noexcept { return {}; }
        void return_value(T x) { result.set_value(std::move(x)); }
        void unhandled_exception() noexcept {
            try { result.set_exception(std::current_exception()); }
            catch (...) { std::terminate(); }
        }
    };
};

FutureTask<int> square_on(Executor& ex, int x) {
    co_await ex.schedule(); // Tu oddajemy sterowanie.
    co_return x * x;        // Wykona sie w watku roboczym.
}

int main() {
    Executor ex;
    auto a = square_on(ex, 6);
    auto b = square_on(ex, 7);
    assert(a.get() + b.get() == 85);
    // Przed zniszczeniem ex oba zadania sie zakonczyly.
}
\end{lstlisting}

\paragraph{Co to porównanie rzeczywiście dowodzi?}
Kod z \texttt{square\_on} jest asynchroniczny względem wywołującego,
ponieważ po \texttt{co\_await} kontynuację przejmuje inny wątek.
Nie oznacza to jeszcze \emph{równoległego wykonywania} obu obliczeń:
executor ma tylko jednego pracownika, więc zadania są przez niego
wykonywane kolejno. Do równoległego przetwarzania dwóch obliczeń
należałoby dodać drugiego pracownika albo wykorzystać istniejącą
pulę wątków.

Pokazany executor jest \emph{szkicem edukacyjnym}. Nie obsługuje
ogólnego anulowania ani ponownego planowania tej samej korutyny,
nie rozwiązuje wszystkich problemów czasu życia i nie stanowi
gotowej biblioteki. W szczególności klient nie ma uchwytu pozwalającego
anulować i bezpiecznie zniszczyć zawieszone zadanie. W systemie
produkcyjnym własność korutyny, zamknięcie executora i czas życia
powiązanych zasobów muszą stanowić jeden, spójny kontrakt.

\section{Obsługa wyjątków, anulowanie i czas życia danych}
\label{sec:cpp15-errors}

\subsection{Wyjątki a \texttt{unhandled\_exception}}

Nieobsłużony wyjątek z ciała korutyny nie może zostać pozostawiony
bez decyzji. Kompilator wywołuje \texttt{promise.unhandled\_exception()}.
W generatorze zapisaliśmy \texttt{std::exception\_ptr}; przy kolejnym
\texttt{next()} błąd zostanie zgłoszony w kodzie klienta. W naszym
\texttt{FutureTask<T>} wyjątek trafia do stanu \texttt{std::future}
i ujawnia się podczas \texttt{get()}. Jest to analogiczne do
propagacji błędów w \texttt{std::async}, ale sam mechanizm wymagał
jawnej implementacji w \texttt{promise\_type}.

Należy też uwzględnić wyjątki zgłaszane przez operacje awaitera
oraz błędy alokacji ramki. Awaria podczas tworzenia ramki może
wystąpić jeszcze przed wykonaniem ciała korutyny. Typ obietnicy
może opcjonalnie udostępniać \texttt{get\_return\_object\_on\_allocation\_failure()}
razem z odpowiednią, niezgłaszającą wyjątku alokacją, ale nie jest to
mechanizm automatycznie dostępny każdemu typowi korutyny.

\subsection{Referencje przechowywane w ramce}

Argumenty referencyjne nie wydłużają czasu życia obiektów.
Jeżeli korutyna przechowuje referencję do zmiennej lokalnej, która
zniknie podczas zawieszenia, jej późniejsze wznowienie może prowadzić
do użycia nieistniejącego już obiektu.

\begin{lstlisting}[caption={Niebezpieczny czas życia referencji w korutynie; fragment koncepcyjny.}]
// Zakladamy lazy Task<int> z poprzedniej sekcji.
Task<int> read_later(const int& ref) {
    co_return ref;
}

// NIEPOPRAWNE:
// auto task = read_later(123); // 123 znika na koncu tej instrukcji
// int value = task.get();     // odczyt przez wiszaca referencje
\end{lstlisting}

Z reguły bezpieczniej przechowywać wewnątrz ramki wartości
przeznaczone do późniejszego użycia (np. obiekt przekazany przez
wartość), a nie referencje do tymczasowych obiektów.
Trzeba też zwracać uwagę na czas życia domknięć lambda:
korutyna będąca lambdą i wykorzystująca przechwycenia może
po zawieszeniu odwoływać się do już zniszczonego obiektu domknięcia.
Sam fakt przeniesienia wynikowej korutyny nie musi temu zapobiegać.

\subsection{Anulowanie jest kooperacyjne}

\texttt{std::stop\_token} (C++20) umożliwia \emph{żądanie}, a nie
wymuszenie zakończenia zadania. Korutyna może sprawdzać token
przed rozpoczęciem kolejnych faz albo korzystać z awaitera,
który potrafi odpowiednio reagować na anulowanie. Nie wynika
z tego jednak, że każde \texttt{co\_await} automatycznie obsługuje
\texttt{stop\_token}.

\begin{lstlisting}[caption={Kooperacyjne anulowanie korutyny: sprawdzenie żądania w punktach kontrolnych; fragment.}]
// Zakladamy Task<int> z poprzedniej sekcji.
#include <stop_token>

Task<int> calculate(std::stop_token token) {
    if (token.stop_requested()) co_return -1;
    // Pierwszy krok obliczen...
    if (token.stop_requested()) co_return -1;
    // Kolejny krok obliczen...
    co_return 42;
}
\end{lstlisting}

W powyższym przykładzie token jest przekazywany przez wartość.
Sam \texttt{stop\_token} nie steruje czasem życia ramki i nie
wznawia korutyny, która czeka na zdarzenie. Do pełnego anulowania
operacji I/O potrzeba jeszcze protokołu z usługą, która wykonuje
rzeczywistą operację, oraz reguł określających \emph{kto} zniszczy
korutynę i \emph{kiedy} wolno to zrobić.

\section{Pułapki synchronizacji i planowania}

\subsection{Równoczesne wznowienie tej samej korutyny}

Nie można traktować \texttt{coroutine\_handle::resume()} jak
operacji atomowej nadającej się do równoczesnego wywoływania przez
wiele wątków. Trzeba zapewnić, że dana zawieszona korutyna ma
jednego logicznego właściciela uprawnienia do wznowienia.
Ponowne wznowienie korutyny już aktywnej lub zakończonej może
prowadzić do niezdefiniowanego zachowania.

\subsection{Utracone powiadomienie i wyścig publikacji uchwytu}

Awaiter może przekazać uchwyt do innego wątku wewnątrz
\texttt{await\_suspend}. Wówczas druga strona może wznowić korutynę
\emph{zanim} pierwotne wywołanie \texttt{await\_suspend} powróci.
Dlatego po opublikowaniu uchwytu kod awaitera nie powinien korzystać
z obiektów, których czas życia mógł się już zakończyć.

W bardziej złożonych awaiterach trzeba również uniknąć zgubienia
zdarzenia kończącego operację między sprawdzeniem gotowości a
zarejestrowaniem oczekiwania. Poprawne rozwiązania stosują np.
mutex, operacje atomowe albo protokół jednorazowego przejęcia
kontynuacji. Samo \texttt{await\_ready()} nie zastępuje synchronizacji.

\subsection{Brak własności \texttt{coroutine\_handle}}

Gdy planista ma w kolejce jedynie \texttt{coroutine\_handle<>},
nie posiada informacji, kto może zniszczyć ramkę. Usunięcie
obiektu \texttt{Task} podczas oczekiwania może prowadzić do
późniejszego wznowienia już zwolnionej pamięci.
Najbezpieczniejszy projekt definiuje własność od początku:
\emph{Task posiada ramkę}, \emph{executor posiada ramkę} albo
\emph{ramka posiada samą siebie}, z dokładnym protokołem przekazania
własności i zakończenia.

\section{Koszty, optymalizacje i pomiary}
\label{sec:cpp15-performance}

Wydajność należy rozpatrywać w całym systemie, a nie tylko przez
porównanie czasu jednego \texttt{co\_await} i jednego \texttt{std::async}.
\begin{itemize}
\item \textbf{Utworzenie wątku.} \texttt{std::async(launch::async, ...)}
      uruchamia pracę w odrębnym wątku; bardzo duża liczba krótkich
      zadań może generować istotny narzut zarządzania wątkami.
\item \textbf{Ramka.} Korutyna może wymagać alokacji pamięci,
      zachowywania stanu i wznawiania. Kompilator może czasem
      wyeliminować alokację ramki, ale nie jest to gwarancja.
\item \textbf{Executor.} Kolejka, mutexy, przełączanie między wątkami
      oraz współdzielenie pamięci również kosztują.
\item \textbf{Lokalność pamięci.} Duża liczba korutyn z rozproszonymi
      ramkami może prowadzić do dodatkowych nietrafień cache.
\item \textbf{Skalowalność.} Dla I/O ważniejsza może być liczba
      oczekujących operacji niż przepustowość pojedynczego zadania.
      Dla CPU-bound kluczowa jest liczba rdzeni, wielkość porcji pracy
      i ograniczenie dostępu do współdzielonych zasobów.
\end{itemize}

Rzetelny benchmark powinien mierzyć co najmniej: całkowitą
przepustowość, opóźnienia (również wartości odstające), liczbę wątków,
liczbę alokacji oraz maksymalne zużycie pamięci. Porównywanie
\texttt{std::async} z \emph{nieuruchamiającą nowych wątków} korutyną
wykonującą natychmiast prostą arytmetykę nie jest miarodajnym
testem wydajności asynchronicznej.

\subsection{Praktyczna decyzja: kiedy używać którego narzędzia?}

\begin{center}
\small
\begin{tabularx}{\textwidth}{@{}lX@{}}
\toprule
\textbf{Problem} & \textbf{Rozsądny pierwszy wybór} \\
\midrule
Jedno niezależne zadanie CPU z wynikiem & \texttt{std::async(std::launch::async, ...)} \\
Wiele krótkich zadań CPU & Pula wątków, executor lub OpenMP \\
Wiele połączeń I/O z długim oczekiwaniem & Korutyny + asynchroniczny runtime I/O \\
Leniwe generowanie sekwencji & Generator korutynowy \\
Wspólny licznik lub krótka sekcja krytyczna & \texttt{atomic} albo mutex, nie korutyna \\
Anulowalne zadanie działające długo & Przemyślany protokół \texttt{stop\_token} \\
Interfejs biblioteki o złożonych zależnościach & \texttt{Task<T>} nad zdefiniowanym executorem \\
\bottomrule
\end{tabularx}
\end{center}

\section{Związek z OpenMP, \texttt{thread} i \texttt{future}}

Rozszerzony rozdział o współbieżności omawiał \texttt{std::thread},
\texttt{std::jthread}, atomiki, blokady, \texttt{std::async} oraz OpenMP.
Korutyny nie unieważniają tych narzędzi; stanowią kolejną warstwę.
\begin{itemize}
\item \texttt{std::thread} opisuje konkretny wątek wykonania,
      a \texttt{std::jthread} dodatkowo pomaga zarządzać jego
      czasem życia i kooperacyjnym zatrzymywaniem.
\item \texttt{std::async} tworzy stan przyszłego wyniku i
      uruchamia wywołanie zgodnie z przyjętą polityką.
\item OpenMP rozdziela pętle i zadania między wątki obsługiwane
      przez środowisko wykonawcze OpenMP; szczególnie przydaje się
      w obliczeniach CPU-bound i kodzie numerycznym.
\item Korutyny pozwalają wygodnie opisać wieloetapową logikę
      oczekiwania i kontynuacji. Ich wydajność i stopień równoległości
      zależą od dostarczonego środowiska wykonawczego.
\end{itemize}

Możliwe jest użycie \emph{korutyn nad pulą wątków}: \texttt{co\_await}
przekazuje kontynuację do kolejki, a pracownik puli później ją wznawia.
Nie należy jednak od razu łączyć wszystkich mechanizmów.
Najpierw określamy, czy program czeka głównie na I/O, czy wykonuje
obliczenia, a następnie wybieramy model wykonania i własności zasobów.

\section{Najczęstsze błędy projektowe}

\begin{enumerate}
\item Zakładanie, że \texttt{co\_await} uruchamia nowy wątek.
\item Próba zwrócenia bezpośrednio \texttt{std::future<T>}
      z własnej korutyny bez odpowiedniego \texttt{promise\_type}.
\item Dwukrotne wywołanie \texttt{destroy()} dla tej samej ramki.
\item Wywołanie \texttt{resume()} po osiągnięciu
      \texttt{final\_suspend}.
\item Pozostawienie referencji lub przechwycenia lambdy do obiektu,
      który znika podczas zawieszenia.
\item Zniszczenie \texttt{Task} mimo tego, że executor posiada
      w kolejce uchwyt do jego ramki.
\item Mylenie automatycznej propagacji wyjątków w \texttt{std::async}
      z obowiązkami implementatora \texttt{promise\_type}.
\item Traktowanie \texttt{std::stop\_token} jako mechanizmu
      przymusowego przerwania dowolnego \texttt{co\_await}.
\item Równoczesne wznawianie tej samej korutyny z kilku wątków.
\item Uznanie, że zawieszenie i wznowienie zawsze kosztuje mniej
      niż zwykłe wywołanie funkcji.
\item Uruchamianie \texttt{std::async} dla tysięcy mikro-zadań
      bez rozważenia puli wątków.
\item Zwracanie z \texttt{await\_suspend} po opublikowaniu uchwytu
      i dalsze używanie potencjalnie zniszczonego awaitera.
\end{enumerate}

\section*{Podsumowanie rozdziału}
\addcontentsline{toc}{section}{Podsumowanie rozdziału}

Korutyny C++20 są mechanizmem \emph{zawieszenia i kontynuacji},
a nie samodzielnym środowiskiem asynchronicznym. Generator używa
\texttt{co\_yield}, aby leniwie udostępniać wartości. Typ zadania
używa \texttt{co\_return}, obiektu \texttt{promise\_type} i awaiterów,
aby opisać wykonanie z wynikiem i wyjątkami. \texttt{co\_await}
przekazuje kontrolę zgodnie z protokołem awaitera, a dopiero executor
lub inne środowisko może wznowić kod w osobnym wątku.

\texttt{std::async} jest prostym narzędziem do uruchomienia funkcji
i odebrania przyszłego wyniku. Korutyny są lepszym budulcem dla
złożonych przepływów wieloetapowych, szczególnie I/O, ale wymagają
większej odpowiedzialności za pamięć, synchronizację, anulowanie
i czas życia ramki. W obu przypadkach należy rozdzielać trzy pytania:
\emph{kto wykonuje kod, kto posiada stan i kto decyduje o zakończeniu}.

\section*{Literatura i dokumentacja}
\addcontentsline{toc}{section}{Literatura i dokumentacja}
\begin{itemize}
\item \url{https://en.cppreference.com/w/cpp/language/coroutines} --- model korutyn C++20.
\item \url{https://en.cppreference.com/w/cpp/coroutine/coroutine_handle} --- uchwyty i zarządzanie ramką.
\item \url{https://en.cppreference.com/w/cpp/coroutine/suspend_always} --- podstawowe awaitery.
\item \url{https://en.cppreference.com/w/cpp/thread/async} --- \texttt{std::async} i polityki uruchamiania.
\item \url{https://en.cppreference.com/w/cpp/thread/future} --- \texttt{std::future}.
\item \url{https://en.cppreference.com/w/cpp/coroutine/generator} --- \texttt{std::generator} (C++23).
\item Lewis Baker, \emph{C++ Coroutines: Understanding the Compiler Transformations}
      --- seria publikacji o modelu wykonania korutyn.
\item David Mazières, \emph{My Tutorial and Take on C++20 Coroutines}, 2021.
\end{itemize}
